% Options for packages loaded elsewhere
\PassOptionsToPackage{unicode}{hyperref}
\PassOptionsToPackage{hyphens}{url}
%
\documentclass[
]{article}
\usepackage{amsmath,amssymb}
\usepackage{iftex}
\ifPDFTeX
  \usepackage[T1]{fontenc}
  \usepackage[utf8]{inputenc}
  \usepackage{textcomp} % provide euro and other symbols
\else % if luatex or xetex
  \usepackage{unicode-math} % this also loads fontspec
  \defaultfontfeatures{Scale=MatchLowercase}
  \defaultfontfeatures[\rmfamily]{Ligatures=TeX,Scale=1}
\fi
\usepackage{lmodern}
\ifPDFTeX\else
  % xetex/luatex font selection
\fi
% Use upquote if available, for straight quotes in verbatim environments
\IfFileExists{upquote.sty}{\usepackage{upquote}}{}
\IfFileExists{microtype.sty}{% use microtype if available
  \usepackage[]{microtype}
  \UseMicrotypeSet[protrusion]{basicmath} % disable protrusion for tt fonts
}{}
\makeatletter
\@ifundefined{KOMAClassName}{% if non-KOMA class
  \IfFileExists{parskip.sty}{%
    \usepackage{parskip}
  }{% else
    \setlength{\parindent}{0pt}
    \setlength{\parskip}{6pt plus 2pt minus 1pt}}
}{% if KOMA class
  \KOMAoptions{parskip=half}}
\makeatother
\usepackage{xcolor}
\usepackage{longtable,booktabs,array}
\usepackage{calc} % for calculating minipage widths
% Correct order of tables after \paragraph or \subparagraph
\usepackage{etoolbox}
\makeatletter
\patchcmd\longtable{\par}{\if@noskipsec\mbox{}\fi\par}{}{}
\makeatother
% Allow footnotes in longtable head/foot
\IfFileExists{footnotehyper.sty}{\usepackage{footnotehyper}}{\usepackage{footnote}}
\makesavenoteenv{longtable}
\setlength{\emergencystretch}{3em} % prevent overfull lines
\providecommand{\tightlist}{%
  \setlength{\itemsep}{0pt}\setlength{\parskip}{0pt}}
\setcounter{secnumdepth}{-\maxdimen} % remove section numbering
\ifLuaTeX
  \usepackage{selnolig}  % disable illegal ligatures
\fi
\IfFileExists{bookmark.sty}{\usepackage{bookmark}}{\usepackage{hyperref}}
\IfFileExists{xurl.sty}{\usepackage{xurl}}{} % add URL line breaks if available
\urlstyle{same}
\hypersetup{
  hidelinks,
  pdfcreator={LaTeX via pandoc}}

\author{}
\date{}

\begin{document}

\hypertarget{maofield-a-dialectical-materialist-framework-for-non-statistical-semantic-representation}{%
\section{MaoField: A Dialectical-Materialist Framework for
Non-Statistical Semantic
Representation}\label{maofield-a-dialectical-materialist-framework-for-non-statistical-semantic-representation}}

\textbf{Author}: Yifan Chen (Chen Yifan) \textbf{Affiliation}:
Independent Researcher \textbf{ORCID}:
\href{https://orcid.org/0009-0008-8344-1149}{0009-0008-8344-1149}
\textbf{Repository}: https://github.com/Wangziqi0/MaoField
\textbf{Concept DOI} (v0.1.0):
\href{https://doi.org/10.5281/zenodo.19550341}{10.5281/zenodo.19550341}
\textbf{Version DOI} (v0.1.1 / this paper):
\href{https://doi.org/10.5281/zenodo.19550342}{10.5281/zenodo.19550342}
\textbf{Date}: 2026-04-20 \textbf{License}: Apache 2.0 (code); CC BY 4.0
(paper)

\begin{center}\rule{0.5\linewidth}{0.5pt}\end{center}

\hypertarget{abstract}{%
\subsection{Abstract}\label{abstract}}

Deep-learning systems trained by gradient descent on statistical
objectives have achieved unprecedented predictive performance, yet
certain qualitatively distinct capacities --- compositional
generalization, structural interpretability, and autonomous semantic
dynamics --- remain out of reach of this paradigm. We propose
\textbf{MaoField}, a candidate \emph{complementary} framework for
problems where structural understanding rather than predictive accuracy
is the operative goal.

The framework's methodology is \textbf{dialectical-materialist}: four
commitments from the Marx--Engels--Lenin--Mao tradition (reflection
theory, theory of contradiction, practice, quantity-to-quality
transition) supply non-optional structural constraints rather than
slogans, and are mapped to concrete design choices (deterministic source
construction, autonomous dynamics, practice-record corpus handling,
regime-transition-aware evaluation).

A \textbf{categorical bridge} to rigorous mathematics is provided by
Lawvere's (1969) identification of adjoint functors with the Hegelian
unity of opposites; we develop the source-attractor map as an
endofunctor \texttt{T\_•\ =\ G∘F\_•} on the category \texttt{Meas} of
measurable spaces, with a conjectured Giry-monad lift \texttt{P∘T\_•} as
the stochastic formalization of synthesis.

MaoField's PDE realization is a Ginzburg-Landau / Allen-Cahn dynamics on
a complex scalar field over a \texttt{32³} lattice. On five BEIR
retrieval benchmarks, the framework improves over a byte-level lexical
baseline by \textbf{+14.7\% to +172.3\%} while lagging SOTA
cross-encoders by 10.8--32.2 pp.~A four-stage diagnostic cascade reveals
that a single numerical attractor count \texttt{k*=2} hides \textbf{two
distinct source-density regimes} (sparse byte / dense BGE), and refines
the open problem of dialectical ascent (OP2) from a single-axis
``non-equilibrium extension needed'' to a \textbf{three-axis co-design
problem} --- source statistics, potential geometry, numerical scheme.

One candidate formalization of Axiom 6 (Banach contraction M2) is
\textbf{falsified}; the axiom itself remains open. We present this as a
\textbf{research-stage position paper} with open problems explicitly
foregrounded, not a completed theory.

\textbf{Keywords}: dialectical materialism, Ginzburg-Landau, category
theory, information retrieval, position paper, Lawvere adjunction, Giry
monad, Kramers rate, spontaneous symmetry breaking, complex scalar
field.

\textbf{arXiv categories}: \texttt{cs.IR} (primary), \texttt{cs.AI},
\texttt{math.CT}, \texttt{cond-mat.dis-nn}.

\begin{center}\rule{0.5\linewidth}{0.5pt}\end{center}

\hypertarget{contents}{%
\subsection{Contents}\label{contents}}

\begin{itemize}
\tightlist
\item
  §1 Introduction
\item
  §2 Categorical Structure

  \begin{itemize}
  \tightlist
  \item
    §2.1 Dialectical-Materialist Foundation
  \item
    §2.2 Adjoint Functors as the Categorical Form of Opposition
  \item
    §2.3 Adjunction, Monad, and the Categorical Structure of Dialectical
    Motion
  \end{itemize}
\item
  §3 Mathematical Formulation
\item
  §4 Experimental Evidence
\item
  §5 Roadmap: Refined Open Problems and Block V Program
\item
  §6 Discussion
\item
  References
\item
  Appendix 2.3.A: Numerical Barrier Computation for A1
\item
  Appendix 3.A: Wirtinger Derivative Convention
\end{itemize}

\begin{center}\rule{0.5\linewidth}{0.5pt}\end{center}

\hypertarget{introduction}{%
\section{1. Introduction}\label{introduction}}

\begin{center}\rule{0.5\linewidth}{0.5pt}\end{center}

\hypertarget{the-paradigm-question}{%
\subsection{1.1 The Paradigm Question}\label{the-paradigm-question}}

Modern artificial intelligence is overwhelmingly built on statistical
pattern-matching at scale: gradient-descent training of high-parameter
neural networks against large data corpora. The empirical successes are
unprecedented, and we do not contest them. We do, however, observe that
the same successes have not produced certain qualitatively distinct
capacities --- robust compositional generalization, structural
transparency, sample-efficient adaptation, and explanatory reasoning
over the field's own outputs --- and that incremental engineering on the
existing paradigm has not closed these gaps. Whether this is a temporary
contingency of architecture, a fundamental limit of statistical
pattern-matching, or a third possibility neither of these
characterizations exhausts is, in our view, an open paradigm question.

This paper takes the third position seriously. \textbf{This is a
position paper}: we articulate a methodological stance and a concrete
candidate framework, with open problems explicitly foregrounded, rather
than asserting full formalization or SOTA performance. We propose a
\textbf{complementary candidate framework} --- not a replacement for
statistical deep learning, but a framework for problems where structural
understanding rather than predictive accuracy is the operative goal. The
candidate is \textbf{field-theoretic, axiom-driven, and rooted in the
methodology of dialectical materialism}. We do not seek SOTA retrieval
performance; our claim is structural interpretability at
useful-above-baseline performance (see §4.2).

\textbf{A note on scope}: the paper's ambition is \textbf{a
methodologically distinct paradigm for artificial intelligence}, not a
contribution to information retrieval per se. Retrieval is chosen as the
\emph{initial empirical domain} because it is \textbf{tractable}:
query--document matching admits a well-posed fixed-point formulation,
benchmarks with established baselines are available, and every step in
the pipeline admits direct mechanistic interpretation. The framework's
intended application scope extends well beyond retrieval --- scientific
hypothesis generation, safety-critical decision systems, educational /
explanatory applications, and more broadly any setting where autonomous
structural dynamics is the operative goal (§6.4, §6.5). Retrieval
results in §4 are evidence that the framework is \textbf{empirically
non-vacuous}, not the ceiling of what the framework is for.

We anticipate two skeptical reactions. First: dialectical materialism,
in 20th-century practice, became a state ideology and is now widely
treated as either museum-piece philosophy or politicized rhetoric ---
neither posture suited to mathematical or empirical work. Second:
applying physics formalism to non-physical systems carries notorious
risks of over-claim. We address both objections directly throughout this
paper. The first by \textbf{using dialectical materialism as a
methodological framework}, not a slogan, and by being explicit about
which axioms are postulated and which are derived (§3.2). The second by
\textbf{explicit attention to the gap between empirical pattern and
rigorous physics} (§3.5, mode-tagging convention), and by foregrounding
open problems rather than concealing them (§5).

\hypertarget{why-a-dialectical-materialist-method}{%
\subsection{1.2 Why a Dialectical-Materialist
Method}\label{why-a-dialectical-materialist-method}}

The dialectical-materialist tradition --- as developed by Marx, Engels,
Lenin, and most pertinently for our purposes in the Maoist contributions
of \emph{On Contradiction} and \emph{On Practice} --- supplies four
methodological commitments that directly inform the framework's design:

\begin{enumerate}
\def\labelenumi{\arabic{enumi}.}
\tightlist
\item
  \textbf{Reflection theory} (反映论, Lenin): cognition is a structural
  correspondence to objective reality, not a statistical fit to
  observation distributions. \emph{Realization in MaoField}: the source
  field \texttt{S(x)} is a deterministic encoding of the input text, not
  a learned compression --- every voxel is interpretable in principle
  (Axiom 4, §3.2)
\item
  \textbf{Theory of contradiction} (矛盾论, Mao): internal opposition is
  the engine of dynamics, not external imposition. \emph{Realization}:
  the field dynamics (Eq. 3.1) are autonomous; no loss function or
  external objective is supplied (Axiom 3)
\item
  \textbf{Theory of practice} (实践论, Mao): cognition deepens through
  practice; the corpus is a record of practice, not a sample from a
  population. \emph{Realization}: the corpus is treated as
  practice-record (Axiom 4)
\item
  \textbf{Quantity-to-quality transition} (质量互变, Engels):
  quantitative changes can produce qualitative shifts that are not
  predictable from extrapolation alone. \emph{Realization}: the
  framework explicitly admits and studies regime transitions (Block IV.5
  Stage A1.4 in §4.9 is a direct experimental instance of this
  principle: a quantitative reduction in barrier height produced a
  qualitative loss of well structure)
\end{enumerate}

These are not mere analogies. Each commitment maps to a concrete
structural choice that distinguishes MaoField from a generic
field-theoretic AI proposal. We conjecture that the convergence of these
commitments produces a framework whose properties are not trivially
replicable by adding any single one of them to a standard pipeline.

\hypertarget{the-categorical-bridge}{%
\subsection{1.3 The Categorical Bridge}\label{the-categorical-bridge}}

A persistent concern with dialectical-materialist approaches has been
the perceived absence of a precise mathematical infrastructure. We adopt
\textbf{Lawvere's 1969 identification of adjoint functors with the
Hegelian unity of opposites} (Lawvere, \emph{Adjointness in
foundations}) as the categorical bridge between dialectical motion and
rigorous mathematics. An adjoint pair \texttt{F\ ⊣\ G} formalizes the
propose-return structure of thesis-antithesis; the induced monad
\texttt{T\ =\ G∘F} (when fully formalized --- see §2.3.3 for current
status) internalizes the synthesis as a T-algebra; the Eilenberg-Moore
and Kleisli categories stratify the modes of dialectical realization.

This categorical infrastructure is \textbf{not} decoration. In §2.3 we
show that it carries conjectural structural content: it identifies the
\textbf{source-density decomposition} (two distinct endofunctors arising
from sparse vs.~dense source constructions, §2.3.2), the
\textbf{reachability gap} as a categorical measure of one of our open
problems (§2.3.4), and the \textbf{three-faces convergence} that unifies
categorical, physical, and phenomenological open problems under a single
conjectured stochastic structure (§2.3.3).

\hypertarget{the-pde-realization-maofield}{%
\subsection{1.4 The PDE Realization:
MaoField}\label{the-pde-realization-maofield}}

The framework's concrete instantiation is a Ginzburg-Landau / Allen-Cahn
dynamics on a complex scalar field \texttt{ψ\ :\ Ω\ →\ ℂ} over a
discrete 3D lattice (32³ voxels):

\begin{verbatim}
γ ∂_t ψ = D ∇²ψ − ∂V/∂ψ* + S(x) + η(x, t)
\end{verbatim}

The functor \texttt{F\ :\ Text\ →\ Field} is the source-construction map
(two implementations studied: byte-frequency \texttt{F\_sparse} and
dense embedding \texttt{F\_dense}); the functor
\texttt{G\ :\ Field\ →\ Score-ready} is the time-evolution to fixed
point. Their composition \texttt{T\ =\ G∘F} is the source-attractor
endofunctor whose categorical structure §2.3 analyzes.

The framework postulates seven axioms (§3.2). Two of these axioms remain
at the level of physical intuition without canonical mathematical
formalization (Axiom 6, OP1) or with refined formulation but unresolved
(Axiom 3, OP2). We treat both as \textbf{open problems prominently
signposted}, not as failures or evasions. A research-stage framework
that conceals its own open problems is a misrepresentation of its
status; we choose the alternative.

\hypertarget{empirical-status-summary}{%
\subsection{1.5 Empirical Status
(Summary)}\label{empirical-status-summary}}

We report experiments on five BEIR retrieval benchmarks (NFCorpus,
SciFact, FiQA, ArguAna, TREC-COVID). Key findings, detailed in §4:

\begin{itemize}
\tightlist
\item
  \textbf{Block I}: MaoField as PDE-based reranker yields
  \textbf{+14.7\% to +172.3\%} improvement over a byte-level lexical
  baseline across all five datasets (Table 4.1, §4.2), while
  underperforming SOTA cross-encoders (BGE-reranker-v2-m3) by
  \texttt{−10.8\ pp} to \texttt{−32.2\ pp}. Additionally, removing
  explicit fusion yields +7.9 to +26.0 pp vs.~MaoField\_base on the same
  datasets, corroborating the Block II finding that fusion erodes
  discriminative signal (§4.3). The framework establishes its empirical
  envelope: meaningfully above naïve baselines, below trained SOTA, with
  structural advantages absent from both
\item
  \textbf{Block II}: Kuramoto diffusive coupling (a candidate for
  replacing explicit fusion) is \textbf{weakly falsified} (max +18.7\%
  gain, below pre-registered 20\% threshold)
\item
  \textbf{Block III}: Attractor enumeration converges on \texttt{k*=2}
  as a \textbf{local silhouette optimum} across three independent
  representations; statistical separation is weak-to-moderate
  (silhouette 0.037--0.223, all below the 0.25 conventional-substantive
  threshold), so \texttt{k*=2} is a candidate structural ceiling subject
  to the methodological caveats in §4.11.2
\item
  \textbf{Block IV.5}: The \texttt{k*\ =\ 2} ceiling decomposes into
  \textbf{two distinct mechanisms} under sparse vs.~dense source fields
  (Stage C, §4.6); pure gradient flow cannot populate distant attractor
  basins (Stage A1, §4.7); isotropic Langevin succeeds for inner
  barriers but fails for outer barriers under finite simulation horizon
  (Stage B1, §4.8; this is a timescale-vs-horizon mismatch rather than a
  paradigm obstruction --- the two-regime framework is formalized in
  §4.10); the OP2 program is refined from ``directed driving needed'' to
  a three-axis co-design problem (Stage A1.4, §4.9)
\end{itemize}

The empirical contribution is \textbf{not} a SOTA performance claim. It
is the \textbf{structural decomposition of a multi-attractor dynamics
problem in a setting where every step admits direct mechanistic
interpretation} --- a property absent from the dominant paradigm.

\hypertarget{open-problems-prominently-foregrounded}{%
\subsection{1.6 Open Problems Prominently
Foregrounded}\label{open-problems-prominently-foregrounded}}

Two open problems organize the framework's relationship to its own
incompleteness:

\begin{itemize}
\tightlist
\item
  \textbf{OP1 (Axiom 6 formalization)}: the principal candidate
  iteration (M2) is \textbf{falsified} by 0-homogeneity (§5.1). Axiom 6
  currently has no known canonical mathematical formalization. We do not
  claim Axiom 6 is wrong; we claim our best candidate for its
  mathematical capture is excluded
\item
  \textbf{OP2 (Axiom 3 non-equilibrium extension)}: the empirical
  refinement from ``needs non-equilibrium'' to ``three-fold co-design
  --- source statistics, potential geometry, numerical scheme'' (§5.2)
  is the central methodological contribution of this paper. Block V
  Phase A (§5.3) is the experimental program designed to address it
\end{itemize}

We foreground OP1 and OP2 because the contributions of §4 and §2.3 make
sense only against a scaffold that includes these problems explicitly
--- and because the \textbf{dialectical-materialist epistemology that
motivates the framework demands that ``open contradictions are the
engine of motion''}, a methodological self-application of Axiom 1.

\hypertarget{contributions}{%
\subsection{1.7 Contributions}\label{contributions}}

Five paper-level contributions. Each is a substantive finding
(categorical, empirical, or methodological-foundational) rather than a
pipeline-level quality-assurance practice; the latter is discussed in
§6.3.

\begin{enumerate}
\def\labelenumi{\arabic{enumi}.}
\item
  \textbf{Categorical structure} for dialectical motion in PDE-based AI:
  Lawvere endofunctor on \texttt{Meas} with conjectured Giry-monad lift
  as the categorical route around deterministic monad obstruction;
  T-algebra reachability subcategory \texttt{T-Alg\^{}\{(η,\ p\_0)\}} as
  the formal measure of OP2; the \textbf{three-faces convergence} of
  categorical / physical / phenomenological open problems onto a single
  conjectured stochastic extension (§2.3)
\item
  \textbf{Source-density mechanism for the \texttt{k*=2} retrieval
  ceiling}: the decomposition reveals \textbf{two structurally distinct
  distributional regimes} under one numerical attractor count --- sparse
  byte = phase concentration vs.~dense BGE = amplitude bistability
  (empirical distributional labels, finite-lattice caveat in §3.5 and
  footnote \footnote{See Section 3.5 (\texttt{{[}\^{}zn-loose{]}}
    footnote) for full text. Summary: \texttt{Z\_n} labels in this paper
    denote \textbf{empirical distributional remaining-symmetry
    patterns}, not rigorous group actions on a thermodynamically-broken
    vacuum. Strict SSB requires (1) group \texttt{G} acting on field
    space, (2) \texttt{V} invariant under \texttt{G}, (3) ground-state
    manifold is \texttt{G}-orbit, (4) non-trivially-transforming order
    parameter \texttt{⟨ψ⟩\ ≠\ 0}, (5) thermodynamic limit
    \texttt{N\ →\ ∞}. Conditions (3)--(5) are not verified on our
    \texttt{N=32³} lattice; \texttt{Z\_n} is a compact descriptor for
    convenience.}). A quantitative by-product central to the result: the
  dense BGE source mean is non-zero at \texttt{t\ =\ −17.1}
  (\texttt{p\ \textless{}\ 10⁻²⁵}), a \textbf{17σ violation of Laplace's
  detailed-balance precondition} that is both the cause of the
  dense-regime bistability and a concrete obstruction to naïve
  composition of learned embeddings with PDE-based frameworks (§4.6,
  §4.9; re-addressed in §6.4)
\item
  \textbf{OP2 refined from single-axis to three-fold co-design}:
  experimental refinement of Axiom 3's open problem from a single-axis
  ``needs non-equilibrium'' formulation to a three-axis (source
  statistics, potential geometry, numerical scheme) problem with
  concrete sub-problem statements and an experimental program. The
  refinement rests on two further findings: \textbf{(i) an
  inner-succeeds/outer-timescale-mismatch reframe} --- isotropic
  Langevin succeeds on the inner barrier (\texttt{ΔV/T\ =\ 16}, observed
  85\% crossing at \texttt{t\_sim\ =\ 250}) but fails on the outer
  barrier (\texttt{ΔV/T\ =\ 105.5}, Kramers escape time exceeds horizon
  by \textbf{\textasciitilde43 orders of magnitude} log-gap),
  establishing that the \texttt{u\ =\ 4} 0\% occupancy is a
  \textbf{timescale--horizon mismatch} rather than a paradigm
  limitation; \textbf{(ii) A1 inner-barrier anomaly} --- the same inner
  crossing exceeds the 1D Kramers prefactor prediction by \textbf{factor
  \textasciitilde10⁴ (\textasciitilde4 orders)}, flagged as an open
  methodological question concerning field-theoretic rate theory (§4.8,
  §4.9, §4.10, §4.11, §5.2)
\item
  \textbf{OP1 M2 falsification}: the principal candidate iteration (M2)
  of Axiom 6's mathematical formalization is \textbf{excluded by
  0-homogeneity obstruction} (§5.1). We do not claim Axiom 6 is wrong;
  we claim our best current candidate for its canonical mathematical
  capture is ruled out, and we signpost the remaining search space
  rather than paper over the failure. This is the framework's most
  explicit instance of an open problem that the research program commits
  to carrying forward
\item
  \textbf{PDE-based retrieval as empirical non-vacuity} (test-bed
  result, not ambition): without retrieval-specific training,
  \texttt{+14.7\%} to \texttt{+172.3\%} improvement over a byte-level
  lexical baseline across five BEIR datasets, \texttt{−10.8\ pp} to
  \texttt{−32.2\ pp} below SOTA cross-encoders, structurally
  interpretable end-to-end. The result establishes that the
  paradigm-level claims of contributions 1--4 are operative on a
  concrete dataset and not purely formal; retrieval is the test bed, not
  the target (§4.2, §6.4)
\end{enumerate}

\hypertarget{outline}{%
\subsection{1.8 Outline}\label{outline}}

\begin{itemize}
\tightlist
\item
  \textbf{§2} \emph{Categorical Structure}: §2.1 dialectical-materialist
  foundation; §2.2 adjoint functors as the categorical form of
  opposition; §2.3 Lawvere endofunctor on \texttt{Meas}, conjectured
  Giry monad lift, three-faces convergence
\item
  \textbf{§3} \emph{Mathematical Formulation}: PDE setup, the seven
  framework axioms with realizations and tension points,
  source-attractor adjunction, OP1/OP2 statements, \textbf{§3.5 on the
  use of symmetry-breaking language under finite-lattice settings}
\item
  \textbf{§4} \emph{Experimental Evidence}: Blocks I-IV (baseline
  retrieval, Kuramoto coupling, attractor enumeration, source-density ×
  operator quadrants); Block IV.5 (Stage C source-density mechanism,
  Stage A1 reachability failure, Stage B1 Langevin and Kramers
  timescale, Stage A1.4 OP2 three-fold diagnostic)
\item
  \textbf{§5} \emph{Roadmap}: §5.1 OP1 M2 falsification; §5.2 OP2
  refined to three-fold co-design; §5.3 Block V Phase A program (A-0
  reachability pre-check, A-1 source statistics, A-2 potential geometry,
  A-3 numerical scheme, A-joint capstone); §5.4 Phase B/C outlook; §5.5
  three-stage publication plan
\item
  \textbf{§6} \emph{Discussion}: limitations, implications, and the
  methodological position of this work in relation to the broader
  dialectical-materialist tradition
\end{itemize}

A companion repository (https://github.com/Wangziqi0/MaoField) provides
the v0.1.0 code and data release (Zenodo concept DOI
\texttt{10.5281/zenodo.19550341}); reproducibility scripts are in
\texttt{paper/reproduce/}.

\hypertarget{categorical-structure}{%
\section{2. Categorical Structure}\label{categorical-structure}}

This section develops the categorical infrastructure that bridges
dialectical-materialist methodology (§2.1) with rigorous mathematics via
Lawvere's adjoint-functor identification (§2.2) and the source-attractor
endofunctor on \texttt{Meas} (§2.3).

\textbf{Position}: §2.1 grounds the dialectical-materialist
methodological commitments in the philosophical tradition; §2.2
introduces the categorical apparatus (adjoint functors) as the bridge
from dialectical motion to rigorous mathematics. §2.3 formalizes the
source-attractor structure as an endofunctor on \texttt{Meas}.

\begin{center}\rule{0.5\linewidth}{0.5pt}\end{center}

\hypertarget{the-dialectical-materialist-methodological-foundation}{%
\subsection{2.1 The Dialectical-Materialist Methodological
Foundation}\label{the-dialectical-materialist-methodological-foundation}}

The methodological commitments stated in §1.2 do not float free of a
tradition. They derive from a coherent body of philosophical work ---
Marx, Engels, Lenin, and the Maoist contributions in \emph{On
Contradiction} (1937) and \emph{On Practice} (1937) --- and they have
specific structural implications for any framework that takes them
seriously. We summarize the four most operative for MaoField, with
explicit forward references to the formalism of §3.

\hypertarget{reflection-theory-ux53cdux6620ux8bba-lenin}{%
\subsubsection{2.1.1 Reflection Theory (反映论,
Lenin)}\label{reflection-theory-ux53cdux6620ux8bba-lenin}}

The Leninist \emph{Materialism and Empirio-Criticism} (1909) defends the
position that cognition is a \textbf{structural correspondence between
subjective representation and objective reality}, and that this
correspondence is achieved through the historical practice of testing
representation against reality. The implication for our framework:
representation should not be a \emph{fit} to a statistical distribution
of observations (which is loss-function-driven and produces compression
artifacts), but a \textbf{deterministic encoding} that preserves the
structure of the input.

\textbf{Realization in MaoField}: the source field \texttt{S(x)} is
constructed deterministically from each input text ---
\texttt{S\_sparse} is byte-frequency scattering (every voxel corresponds
to a specific UTF-8 byte value; position information is collapsed to
frequency, so the construction is injective at the histogram level but
lossy at the sequence level); \texttt{S\_dense} is BGE-M3 embedding
tiled (every channel directly traceable to an embedding component, even
if the embedding itself was learned upstream). We do not claim
sequence-level positional reconstruction; we claim every voxel value is
interpretable with respect to the input. This is the structural
realization of Reflection Theory at the encoding boundary.

\hypertarget{theory-of-contradiction-ux77dbux76feux8bba-mao}{%
\subsubsection{2.1.2 Theory of Contradiction (矛盾论,
Mao)}\label{theory-of-contradiction-ux77dbux76feux8bba-mao}}

Mao's 1937 \emph{On Contradiction} identifies the \textbf{internal
opposition within a system as the source of its motion}. External
factors are conditions; internal contradictions are causes. This implies
that a dynamical system's evolution should be driven by structure
intrinsic to the system, not by an externally imposed objective
function.

\textbf{Realization in MaoField}: the dynamics (Eq. 3.1) are autonomous
in the strict sense --- the system contains no loss function, no
training signal, no externally imposed gradient. The drives are: (i) the
source field \texttt{S} (encoding input regularity), (ii) the potential
\texttt{V} (encoding intrinsic field structure), and (iii) under
non-equilibrium extension, the noise term \texttt{η} (representing
thermal coupling, though see §3.2 Tension Point 2). The convergence to a
fixed point is the working-out of internal contradictions among
\texttt{S}, \texttt{V}, and the field's own configuration --- not
optimization toward a specified target.

\hypertarget{theory-of-practice-ux5b9eux8df5ux8bba-mao}{%
\subsubsection{2.1.3 Theory of Practice (实践论,
Mao)}\label{theory-of-practice-ux5b9eux8df5ux8bba-mao}}

Mao's 1937 \emph{On Practice} posits that \textbf{cognition deepens
through repeated practice}, and that knowledge accumulates not as a
static accumulation but as an iterative refinement of the
practice-cognition cycle. This implies that the corpus on which a
framework operates is not a \emph{sample from a population} (the framing
implicit in standard statistical learning) but a \textbf{record of
practice} --- and the framework's job is to extract structural
invariants from this record, not to fit a generative model of how the
record was sampled.

\textbf{Realization in MaoField}: the corpus is treated as
practice-record (Axiom 4, §3.2). The framework's evaluation on retrieval
benchmarks (Block I, §4.2) measures structural invariants (basin
assignment consistency, attractor enumeration) rather than generative
likelihood. We do not train an upstream model on the BEIR corpus; we
treat each query-document pair as a single instance of the practice of
matching, and the dynamics' fixed point as the structural outcome of
that practice.

\hypertarget{quantity-to-quality-transition-ux8d28ux91cfux4e92ux53d8-engels}{%
\subsubsection{2.1.4 Quantity-to-Quality Transition (质量互变,
Engels)}\label{quantity-to-quality-transition-ux8d28ux91cfux4e92ux53d8-engels}}

Engels's \emph{Dialectics of Nature} (1873--1883, posthumously published
1925) and the subsequent Marxist tradition develop the principle that
\textbf{quantitative changes can produce qualitative shifts that are not
predictable from extrapolation alone}. The implication for system
design: regime transitions must be admitted as a structural possibility,
not glossed over by smoothing.

\textbf{Realization in MaoField}: the framework explicitly admits and
studies regime transitions. The single most direct experimental instance
is Stage A1.4 (§4.9): an engineered quantitative reduction of the
potential's outer barrier --- undertaken to probe whether continuity of
parameter change would yield continuity of behavior --- instead
triggered a qualitative regime shift in which the field explored
amplitudes orders of magnitude larger than the original potential
admitted. §4.9.3 decomposes this regime shift into three coupled axes
(source statistics, potential geometry, numerical scheme); the
quantitative-qualitative discontinuity is the phenomenon, and the
three-axis decomposition is its mechanism. The barrier reduction was not
a continuous-improvement engineering attempt --- it was an experimental
probe of the regime-transition principle, and its outcome (the
three-fold co-design refinement of §5.2) is precisely the kind of
qualitative restructuring that quantity-to-quality transition predicts.

\hypertarget{three-methodological-implications}{%
\subsubsection{2.1.5 Three Methodological
Implications}\label{three-methodological-implications}}

The four commitments above jointly imply three positions that
distinguish MaoField from a generic field-theoretic AI proposal:

\begin{enumerate}
\def\labelenumi{\arabic{enumi}.}
\tightlist
\item
  \textbf{No statistical fitting at the framework boundary}: source
  construction is deterministic; there is no end-to-end gradient that
  adjusts the framework's parameters against an objective. Statistical
  methods may appear \emph{internally} (e.g., BGE embeddings are
  learned) but only as inputs that the framework subsequently subjects
  to its own non-statistical dynamics
\item
  \textbf{Open problems are constitutive of the research program}: the
  Theory of Contradiction implies that a closed framework --- one with
  no internal contradictions --- would be one with no source of motion.
  We expose OP1 and OP2 prominently (§5) because their persistence is
  the framework's evidence of life
\item
  \textbf{Regime transitions are first-class objects}: Stage A1.4 (§4.9)
  is not an ``experimental error'' or a ``failed attempt''; it is the
  framework's own demonstration of the quantity-to-quality principle in
  action. The three-fold co-design refinement of OP2 is the framework's
  response to this transition
\end{enumerate}

These three positions --- together with Lawvere's categorical bridge
(§2.2) --- define the operational space within which MaoField's specific
choices (PDE form, axioms, source constructions) can be evaluated.

\begin{center}\rule{0.5\linewidth}{0.5pt}\end{center}

\hypertarget{adjoint-functors-as-the-categorical-form-of-opposition}{%
\subsection{2.2 Adjoint Functors as the Categorical Form of
Opposition}\label{adjoint-functors-as-the-categorical-form-of-opposition}}

A persistent obstacle to the mathematical respectability of dialectical
methods has been the perceived absence of a formal apparatus that
captures the structure of opposition without collapsing into either
idealist abstraction or rhetorical analogy. This obstacle was
substantially addressed by \textbf{Lawvere's 1969 identification of
adjoint functors with the Hegelian unity of opposites}
(\emph{Adjointness in foundations}, Dialectica 23:281--296). We adopt
this identification as the categorical foundation of our framework.

\hypertarget{adjunctions-recall}{%
\subsubsection{2.2.1 Adjunctions Recall}\label{adjunctions-recall}}

For categories \texttt{C} and \texttt{D}, an \textbf{adjunction}
\texttt{F\ ⊣\ G} consists of functors \texttt{F\ :\ C\ →\ D} and
\texttt{G\ :\ D\ →\ C} together with a natural bijection

\begin{verbatim}
Hom_D(F(c), d) ≅ Hom_C(c, G(d))                                       (2.1)
\end{verbatim}

for all objects \texttt{c\ ∈\ C} and \texttt{d\ ∈\ D}. Equivalently, an
adjunction is specified by two natural transformations --- the
\textbf{unit} \texttt{η\ :\ Id\_C\ ⇒\ G∘F} and the \textbf{counit}
\texttt{ε\ :\ F∘G\ ⇒\ Id\_D} --- satisfying the \textbf{triangle
identities}:

\begin{verbatim}
(εF) ∘ (Fη) = id_F          and          (Gε) ∘ (ηG) = id_G            (2.2)
\end{verbatim}

The triangle identities encode the \textbf{closure} of the
propose-return cycle: composing the propose (\texttt{F}) with the unit
(\texttt{η}) and the counit (\texttt{ε}) recovers \texttt{F};
symmetrically for \texttt{G}.

\hypertarget{lawveres-identification}{%
\subsubsection{2.2.2 Lawvere's
Identification}\label{lawveres-identification}}

Lawvere's contribution was to identify the structure (2.1)--(2.2) with
the Hegelian dialectical structure of \emph{unity of opposites}:

\begin{itemize}
\tightlist
\item
  \texttt{F} is the \textbf{propose} functor (the moment of
  \emph{thesis}, ``立''): it sends each object of \texttt{C} to its
  image in \texttt{D}. In dialectical-materialist language, \texttt{F}
  is the practice that actively constitutes a representation
\item
  \texttt{G} is the \textbf{return} functor (the moment of
  \emph{antithesis}, ``破''): it sends each object of \texttt{D} back to
  \texttt{C}. In dialectical-materialist language, \texttt{G} is the
  reflection that re-grounds the proposed representation in its source
  category
\item
  The \textbf{unit} \texttt{η\ :\ Id\_C\ ⇒\ G∘F} captures the
  \textbf{non-trivial return}: an object \texttt{c\ ∈\ C}, after being
  proposed into \texttt{D} and returned to \texttt{C}, is in general
  \emph{not} identical to \texttt{c}. The discrepancy
  \texttt{c\ →\ (G∘F)(c)} encodes how the propose-return cycle changes
  the original
\item
  The \textbf{counit} \texttt{ε\ :\ F∘G\ ⇒\ Id\_D} symmetrically
  captures the discrepancy on the \texttt{D} side
\item
  The \textbf{triangle identities} encode the \textbf{closure of the
  dialectical cycle}: the propose-return-propose composition recovers
  the original propose, and symmetrically. This closure is the
  categorical form of dialectical determinism --- the cycle is internal,
  not arbitrary
\end{itemize}

This identification is \textbf{rigorous}: each element of the
dialectical structure (thesis, antithesis, closure) maps to a precise
categorical element (functor, natural transformation, identity). It is
also \textbf{non-trivial}: not every adjunction is a useful dialectical
structure (most are degenerate; the structure becomes substantive only
when \texttt{F}, \texttt{G} are non-isomorphism functors), and the
dialectical reading constrains which adjunctions are physically
meaningful.

\hypertarget{the-monad-as-internalized-dialectic}{%
\subsubsection{2.2.3 The Monad as Internalized
Dialectic}\label{the-monad-as-internalized-dialectic}}

The composite \texttt{T\ :=\ G∘F\ :\ C\ →\ C} is, when fully formalized
as a \textbf{monad}, the internalization of the entire dialectical cycle
within \texttt{C}. The monad multiplication
\texttt{μ\ :=\ G·ε·F\ :\ T²\ ⇒\ T} collapses two iterations of the
propose-return cycle into one, formalizing the \emph{iterative
deepening} of dialectical motion. A \textbf{T-algebra}
\texttt{(X,\ α:\ TX\ →\ X)} is an object of \texttt{C} that is fixed
under the iterated cycle --- the categorical form of \textbf{synthesis}
(Aufhebung), in which the propose-return motion has stabilized into a
self-consistent structure.

§2.3 develops this construction for MaoField's source-attractor
dynamics, identifying:

\begin{itemize}
\tightlist
\item
  the explicit functors \texttt{F\_sparse,\ F\_dense,\ G} for the byte
  and dense source constructions,
\item
  the conjectured Giry-monad lift required to upgrade the deterministic
  endofunctor \texttt{T\_•} to a genuine monad (§2.3.3),
\item
  the \textbf{reachability subcategory} \texttt{T-Alg\^{}\{(η,\ p\_0)\}}
  that distinguishes existent from realizable dialectical structures
  (§2.3.4),
\item
  and the \textbf{three-faces convergence} that links our categorical,
  physical, and phenomenological open problems under a single
  conjectured stochastic structure (§2.3.3).
\end{itemize}

\hypertarget{what-the-categorical-bridge-buys-us}{%
\subsubsection{2.2.4 What the Categorical Bridge Buys
Us}\label{what-the-categorical-bridge-buys-us}}

We close §2.2 with an explicit statement of what the categorical bridge
accomplishes --- and what it does not.

\textbf{It does provide}: a precise vocabulary for describing
dialectical motion (\texttt{F\ ⊣\ G}, \texttt{η}, \texttt{ε},
\texttt{T}-algebra) that is shared with other branches of mathematics; a
language in which the framework's structural commitments can be checked
against well-understood categorical theorems (e.g., the existence of
T-algebras under the Eilenberg-Moore construction, the compositional
behavior of adjunctions); and a setting in which open problems can be
precisely formulated (§2.3.4's \texttt{Δ\_OP2} as cardinality complement
is one such formulation).

\textbf{It does not provide}: a guarantee that the
dialectical-materialist commitments of §2.1 are uniquely realized by the
categorical formalism (other categorical structures may admit
dialectical readings); a substitute for empirical work (the categorical
structure must be \emph{realized} in a concrete dynamics, and §3 onward
make this realization explicit); or a complete formalization of the
framework's open problems (§2.3.3 explicitly notes that the monad
structure is conjectured, not proven, in the deterministic case).

The categorical bridge is \textbf{a methodological instrument}, not a
foundational guarantee. In §2.3 we show that it is sharp enough to
identify structural decompositions (source-density), reachability gaps
(\texttt{Δ\_OP2}), and convergences across apparently independent open
problems (three-faces) --- these are the deliverables we claim. The
bridge is \textbf{not} a proof that MaoField is the unique categorical
realization of dialectical materialism, and we do not claim it is.

\begin{center}\rule{0.5\linewidth}{0.5pt}\end{center}

\emph{Forward reference}. §2.3 below develops the explicit
monad-theoretic structure of MaoField's source-attractor dynamics,
including the source-density decomposition into two endofunctors
\texttt{T\_sparse} and \texttt{T\_dense} (§2.3.2), the explanation of
why deterministic gradient flow does not admit a monad structure and
what stochastic extension is required (§2.3.3), and the categorical
formalization of OP2 as a reachability gap in the T-algebra category
(§2.3.4).

\hypertarget{adjunction-monad-and-the-categorical-structure-of-dialectical-motion}{%
\subsection{2.3 Adjunction, Monad, and the Categorical Structure of
Dialectical
Motion}\label{adjunction-monad-and-the-categorical-structure-of-dialectical-motion}}

\textbf{Position}: position-paper-level identification. The bare
endofunctor picture on \texttt{Meas} is rigorous; the monad lift via
Giry is conjectured and reserved for follow-up work. All \texttt{Z\_n}
labels follow the loose-sense convention defined in §3.5 and footnote
\texttt{{[}\^{}zn-loose{]}}.

\begin{center}\rule{0.5\linewidth}{0.5pt}\end{center}

\hypertarget{from-adjunction-to-monad-the-closure-of-dialectical-operation}{%
\subsection{2.3.1 From Adjunction to Monad: The Closure of Dialectical
Operation}\label{from-adjunction-to-monad-the-closure-of-dialectical-operation}}

An adjoint pair \texttt{F\ ⊣\ G} between categories \texttt{C} and
\texttt{D} expresses a primitive form of dialectical opposition:
\texttt{F} \emph{proposes} objects of \texttt{C} into \texttt{D} (the
moment of \textbf{thesis}, ``立''), while \texttt{G} \emph{returns}
objects of \texttt{D} back into \texttt{C} (the moment of
\textbf{antithesis}, ``破''). The adjunction is constituted by two
natural transformations --- the \textbf{unit} \texttt{η:\ Id\_C\ ⇒\ G∘F}
and the \textbf{counit} \texttt{ε:\ F∘G\ ⇒\ Id\_D} --- satisfying the
triangle identities, which encode the \textbf{closure} of the
propose-return cycle.

The composition \texttt{T\ :=\ G∘F\ :\ C\ →\ C} is, when fully
formalized, a \textbf{monad} with multiplication
\texttt{μ\ :=\ G·ε·F\ :\ T²\ ⇒\ T}. Whereas the adjunction alone
describes a single dialectical oscillation, the monad \texttt{T}
internalizes the \emph{iterative} dialectical process entirely within
\texttt{C}. A \textbf{T-algebra} \texttt{(X,\ α:\ TX\ →\ X)} is
precisely an object that has reached a self-consistent state under this
iteration --- it is the categorical form of the
\textbf{synthesis}\footnote{\emph{Synthesis} (Aufhebung in Hegelian
  terminology, capturing the dual sense of cancellation and
  preservation; often mistranslated as ``combination'', which loses the
  negation-of-negation structure). The T-algebra formalization preserves
  this dual structure: the algebra map \texttt{α:\ TX\ →\ X} both
  \emph{absorbs} (cancels) the monadic iteration and \emph{retains}
  (preserves) its action as an endomorphism.}: not a third point
adjoined, but a structural invariant of the propose-return-reintegrate
cycle. The Kleisli and Eilenberg-Moore categories of \texttt{T} stratify
the algebras by their mode of realization.

This framing is due to Lawvere (1969); our contribution is to
\textbf{propose} a physical realization of the structure under PDE
dynamics and to \textbf{identify} two subtleties that pure categorical
abstraction elides: the \textbf{decomposition of the operator F by
source density} (§2.3.2), and the \textbf{failure of deterministic
gradient flow to satisfy monad axioms} (§2.3.3). We do not claim our
realization is the unique or the canonical one; we claim it is rigorous
at the endofunctor level and that its open problems organize a coherent
research program.

\hypertarget{two-paths-two-endofunctors-the-source-density-decomposition}{%
\subsection{2.3.2 Two Paths, Two Endofunctors: The Source-Density
Decomposition}\label{two-paths-two-endofunctors-the-source-density-decomposition}}

MaoField implements the functor \texttt{F} via a source-field
construction \texttt{Text\ →\ Field\ ∈\ ℂ\^{}N}. Our experiments reveal
that this construction \textbf{is not unique}, and that the specific
realization determines the algebraic structure.

Let \texttt{F\_sparse\ :\ Text\ →\ Field} denote the byte-frequency
construction (UTF-8 bits scattered over the 32³ grid,
\textasciitilde10\% occupancy), and let
\texttt{F\_dense\ :\ Text\ →\ Field} denote the dense construction via a
pretrained embedding (here BGE-M3, tiled to full grid coverage). With a
common \texttt{G\ :\ Field\ →\ Score-ready\ representation}
(Ginzburg-Landau gradient flow to fixed point), the composition
\texttt{G∘F\_•} properly spans three distinct layers (Text, Field,
Score-ready). To discuss categorical structure, we embed all three as
subcategories of an ambient category \texttt{Meas} of measurable spaces,
extending \texttt{F\_•,\ G} by identity outside their natural domains so
that \texttt{G∘F\_•\ :\ Meas\ →\ Meas} becomes a well-defined
endofunctor. Under this embedding we obtain two composite
\textbf{endofunctors} (the lift to genuine monads is reserved for §2.3.3
and footnote \footnote{Concretely, \texttt{Meas} is the category of
  measurable spaces with measurable maps; Text, Field, and Score-ready
  objects are embedded via their natural σ-algebras (discrete for Text,
  Borel for the ℂ\^{}N Field, Borel for Score). Deterministic maps
  \texttt{F\_•,\ G} extend to identity on the complement, so
  \texttt{T\_•\ :=\ G∘F\_•} is a bona fide endofunctor of \texttt{Meas}.
  The Giry monad \texttt{P} is the standard probability-measure monad on
  \texttt{Meas} (Giry 1982; Jacobs 2010): for measurable space
  \texttt{(X,\ Σ\_X)}, \texttt{P(X)} is the space of probability
  measures on \texttt{(X,\ Σ\_X)} with the σ-algebra generated by
  evaluation maps \texttt{μ\ ↦\ μ(A)} for \texttt{A\ ∈\ Σ\_X}; the unit
  \texttt{η\_X\ :\ X\ →\ P(X)} is the Dirac map \texttt{x\ ↦\ δ\_x}; the
  multiplication \texttt{μ\_X\ :\ P(P(X))\ →\ P(X)} is the
  marginalization \texttt{Π\ ↦\ ∫ν\ dΠ(ν)}. The conjectured lift
  \texttt{P\ ∘\ T\_•} of our endofunctor would carry the (η, μ) of
  \texttt{P} provided the natural-transformation diagrams compose;
  verifying or disproving this is OP1's category-theoretic component,
  related to but distinct from the M2-falsification of OP1's analytic
  component (see §5.1).}):

\begin{itemize}
\tightlist
\item
  \texttt{T\_sparse\ :=\ G∘F\_sparse}: phase concentrated
  (\texttt{\textbar{}R\textbar{}\_per-doc\ =\ 0.84}), U(1) broken in the
  empirical distribution but not antipodally split --- \texttt{Z\_2} not
  realized in the phase sector. Phase concentration to a narrow band
  (kmeans-2 center distance 26.2°). See §4.6 for full diagnostics.
\item
  \texttt{T\_dense\ :=\ G∘F\_dense}: amplitude bimodal at
  \texttt{\textbar{}ψ\textbar{}∈\{0,1\}}, phase much less concentrated
  (\texttt{\textbar{}R\textbar{}\_per-doc\ =\ 0.20}, still non-uniform
  vs uniform-baseline 0.006). \texttt{V=¼(\textbar{}ψ\textbar{}²−1)²}
  admits only \texttt{\textbar{}ψ\textbar{}=1} as a true minimum;
  \texttt{\textbar{}ψ\textbar{}=0} is an unstable fixed point with
  \texttt{V(0)=¼\ ≠\ V(1)=0}, so no \texttt{Z\_2} group action
  interchanges the two basins. The ``amplitude bistability'' is a
  \textbf{distributional pattern, not a group-theoretic claim}.
\end{itemize}

Numerically both yield \texttt{k*=2}. \textbf{Categorically they
correspond to distinct endofunctors}: \texttt{T\_sparse}'s candidate
Kleisli category is phase-degenerate; \texttt{T\_dense}'s is
amplitude-stratified.\footnote{This is a \emph{conjectured} categorical
  interpretation based on the experimental phase/amplitude observations;
  a fully rigorous functorial treatment (e.g., explicit functors to
  \texttt{Z\_n}-equivariant categories indexed by the broken symmetry
  pattern) is left to future work.} The claim that ``dialectical closure
is determined by the operator alone'' is refuted by this decomposition
--- \textbf{the source-field functor \texttt{F} co-determines the
algebra structure}. This is a categorical formulation of the materialist
thesis that \emph{practice (the mode of source-field construction)
conditions the forms of cognitive closure}.

\hypertarget{why-deterministic-gradient-flow-is-not-a-monad-and-what-this-reveals}{%
\subsection{2.3.3 Why Deterministic Gradient Flow Is Not a Monad --- and
What This
Reveals}\label{why-deterministic-gradient-flow-is-not-a-monad-and-what-this-reveals}}

A subtlety arises when one attempts to upgrade the endofunctor
\texttt{T\_•\ :=\ G∘F\_•} to a genuine monad in the Lawvere sense. The
composition \texttt{G∘F\_•} is an endofunctor of \texttt{Meas} once
embedded as in §2.3.2, but the multiplication \texttt{μ\ :\ T²\ ⇒\ T}
requires that the iterative composition \texttt{T∘T} admit a canonical
natural transformation back to \texttt{T}. For \textbf{deterministic}
dynamics --- gradient flow \texttt{∂\_t\ ψ\ =\ −δH/δψ*} --- this is
\textbf{not the case}.

\hypertarget{the-categorical-diagnosis}{%
\subsubsection{The Categorical
Diagnosis}\label{the-categorical-diagnosis}}

Deterministic gradient flow is naturally an \texttt{R⁺}-monoid action on
field space --- equivalently, an \textbf{F-coalgebra} for the
time-translation endofunctor (\texttt{X\ →\ F(X)} describing how state
evolves under one time-step). It is not a monad: there is no canonical
multiplication \texttt{μ:\ F∘F\ →\ F} satisfying the monad coherence
axioms (associativity and unit) in the deterministic regime. The
standard categorical lift to a genuine monad goes through
\textbf{stochastic extension}: replacing deterministic flows with Markov
kernels yields the \textbf{Giry monad} \texttt{P} on \texttt{Meas} (Giry
1982; Jacobs 2010). We conjecture that the source-attractor endofunctor
\texttt{T\_•} admits a Giry-monad lift \texttt{P\ ∘\ T\_•} inheriting
(η, μ) from the underlying stochastic dynamics. Rigorous formalization
of \texttt{P\ ∘\ T\_•} is an open problem.

\hypertarget{the-physical-diagnosis}{%
\subsubsection{The Physical Diagnosis}\label{the-physical-diagnosis}}

The same stochastic extension required for the categorical monad lift is
\textbf{also the physical mechanism} required to overcome the Kramers
timescale obstruction observed in Stage B1 (§4.8). Under deterministic
gradient flow alone, our system cannot escape any single basin of
attraction (\texttt{F} is monotone non-increasing along trajectories);
under isotropic Langevin, escape is possible but the Kramers rate
\texttt{∝\ exp(−ΔV/T\_eff)} makes high-barrier crossings unreachable
within finite simulation horizons. \textbf{Closing OP2 (Axiom 3's open
problem of non-equilibrium extensions) requires moving beyond pure
gradient flow to genuinely stochastic / non-equilibrium dynamics} ---
which is the same operation that supplies the Giry-monad multiplication
\texttt{μ} for the categorical lift.

\hypertarget{three-faces-of-one-open-problem}{%
\subsubsection{Three Faces of One Open
Problem}\label{three-faces-of-one-open-problem}}

The categorical structure (§2.3.1--2.3.2), the OP2 physical roadmap
(§5.2), and the Stage B1 phenomenology (§4.8) are not three independent
open problems --- they are \textbf{three faces of a single conjectured
stochastic extension}:

\begin{verbatim}
  Categorical face        Physical face          Phenomenological face
  ───────────────         ─────────────          ─────────────────────
  Endofunctor T_•   ⟶    Gradient flow:    ⟶   B1: Inner barrier 
  not a monad             monotone, no            crossed (Kramers
  in deterministic        escape from any         marginal); outer
  regime                  basin                   barrier at 43-orders-
                                                  of-magnitude gap
                                                  
                          ┃                       ┃
                          ┃                       ┃
                          ┃ requires same         ┃
                          ┃ stochastic            ┃
                          ┃ extension             ┃
                          ┃                       ┃
                          ▼                       ▼

               ╔══════════════════════════════════════════╗
               ║ Conjectured stochastic structure:         ║
               ║                                           ║
               ║  Categorical:  Giry monad lift P ∘ T_•   ║
               ║  Physical:     directed non-equilibrium  ║
               ║                driving (OP2 path b)       ║
               ║  Phenom.:      barrier-and-confinement   ║
               ║                co-design (OP2 path a)     ║
               ╚══════════════════════════════════════════╝
\end{verbatim}

This \textbf{convergence of three apparently independent open problems}
onto a single conjectured stochastic structure is, in our reading, the
\textbf{most coherent organizing conjecture currently visible} for
MaoField as a research program. The dialectical-materialist principle
that ``categorical formalization, physical mechanism, and empirical
observation should converge'' (compatible with Axiom 4's
materially-grounded epistemology) is \textbf{suggested}, not proven, by
the framework's open problems lining up this way; rigorous demonstration
of the convergence is itself part of the open problem.

\hypertarget{reachability-the-categorical-measure-of-open-problem-2}{%
\subsection{2.3.4 Reachability: The Categorical Measure of Open Problem
2}\label{reachability-the-categorical-measure-of-open-problem-2}}

The source-density decomposition in §2.3.2 locates the choice of algebra
structure at the functor \texttt{F}. A second, independent axis of
constraint emerges from the \textbf{dynamics}: even when the target
algebraic structure is in-principle present in \texttt{T-Alg\_T} (the
category of T-algebras under the conjectured monad lift), the dynamics
of MaoField may fail to reach it. Stage A1 of our diagnostic cascade
(§4.7) makes this concrete.

When \texttt{V(ψ)} is extended from the double-well
\texttt{¼(\textbar{}ψ\textbar{}²−1)²} to a triple-well
\texttt{\textbar{}ψ\textbar{}²(\textbar{}ψ\textbar{}²−1)²(\textbar{}ψ\textbar{}²−4)²},
the would-be T-algebra category gains additional objects corresponding
to the \texttt{\textbar{}ψ\textbar{}=2} basin. However, under
gradient-flow dynamics initialized from
\texttt{\textbar{}ψ\textbar{}\_init\ ≈\ 1.0\ ±\ 0.3}, the observed
occupancy of the \texttt{\textbar{}ψ\textbar{}=2} basin is
\textbf{0.00\%} (99.15\% at \texttt{\textbar{}ψ\textbar{}=1}, 0.85\% at
\texttt{\textbar{}ψ\textbar{}=0}). The \texttt{u=4} algebras exist in
the theory but are not realized in the experiment.

This motivates the following refinement. Define the \textbf{reachable
subcategory} under a unit \texttt{η} and an initial distribution
\texttt{p\_0}:

\begin{quote}
\texttt{T-Alg\_T\^{}\{(η,\ p\_0)\}\ :=\ \{\ (X,\ α)\ ∈\ T-Alg\_T\ \textbar{}\ ∃\ x\_0\ ∈\ supp(p\_0),\ η\_\{x\_0\}\ generates\ a\ Kleisli-reachable\ path\ to\ (X,\ α)\ via\ μ\ \}}
\end{quote}

In general \texttt{T-Alg\_T\^{}\{(η,\ p\_0)\}\ ⊊\ T-Alg\_T}. We propose
the \textbf{categorical gap}

\begin{quote}
\texttt{Δ\_\{OP2\}\ :=\ \textbar{}T-Alg\_T\textbar{}\ ⊖\ \textbar{}T-Alg\_T\^{}\{(η,\ p\_0)\}\textbar{}}
\end{quote}

as the \textbf{formal measure of Open Problem 2}: the unrealized
algebraic potential of dialectical structures that are
\emph{mathematically present} in the (conjectured) monad but
\emph{dynamically unreachable} under descending-only (gradient-flow)
evolution.\footnote{\texttt{Δ\_\{OP2\}} as defined above uses
  cardinality complement as a symbolic first pass; alternative
  quantifications --- categorical entropy, persistent-homology dimension
  of the unreached manifold (Adams et al., giotto-tda 2021), Kan
  extension obstructions --- remain follow-up work. The qualitative
  invariant --- the gap between \emph{existent} and \emph{realizable}
  T-algebras --- is conjectured (not proven) to be stable across these
  choices. Closing this gap is one component of OP2 (§5.2); the other
  component is the empirical co-design refinement identified in §4.9
  (A1.4 verdict).}

Two technical caveats are required for \texttt{Δ\_\{OP2\}} to be
well-defined:

\begin{enumerate}
\def\labelenumi{(\roman{enumi})}
\item
  \textbf{The notation \texttt{⊖}} denotes full-subcategory complement:
  \texttt{Δ\_\{OP2\}} is intended as the (essentially small) full
  subcategory of \texttt{T-Alg\_T} whose objects do \textbf{not} lie in
  \texttt{T-Alg\_T\^{}\{(η,\ p\_0)\}}, equipped with the inherited
  morphisms. \texttt{\textbar{}·\textbar{}} is then the cardinality (or,
  for proper-class issues, the cardinality after restriction to a
  Grothendieck universe), or --- preferably --- a coarser invariant such
  as groupoid cardinality \texttt{Σ\ 1/\textbar{}Aut\textbar{}} over
  isomorphism classes.
\item
  \textbf{\texttt{T-Alg\_T} may be a proper class} in the absence of a
  small-set ambient universe; we restrict attention to
  \texttt{T-Alg\_T\ ∩\ U} for a suitable universe \texttt{U} containing
  all algebra structures of finite presentation under the conjectured
  monad lift, and treat the cardinality of the complement within this
  restricted setting. Footnote \footnote{\texttt{Δ\_\{OP2\}} as defined
    above uses cardinality complement as a symbolic first pass;
    alternative quantifications --- categorical entropy,
    persistent-homology dimension of the unreached manifold (Adams et
    al., giotto-tda 2021), Kan extension obstructions --- remain
    follow-up work. The qualitative invariant --- the gap between
    \emph{existent} and \emph{realizable} T-algebras --- is conjectured
    (not proven) to be stable across these choices. Closing this gap is
    one component of OP2 (§5.2); the other component is the empirical
    co-design refinement identified in §4.9 (A1.4 verdict).}
  acknowledges that other quantifications (categorical entropy,
  persistent-homology dimension of the unreached manifold, Kan extension
  obstructions) are equally valid and conjecturally agree in a common
  limit; the cardinality complement is a symbolic first pass.
\end{enumerate}

The ``Kleisli-reachable path'' condition in the definition of
\texttt{T-Alg\_T\^{}\{(η,\ p\_0)\}} should be read as: there exists a
finite composition of Kleisli morphisms (under the conjectured monad
lift; see §2.3.3) starting from a unit of an element in
\texttt{supp(p\_0)} and reaching the algebra \texttt{(X,\ α)} as an
object. This is a categorical translation of the dynamical reachability
condition; rigorous formulation requires the Kleisli composition
\texttt{g\ ∘\_K\ f\ :=\ μ\ ∘\ Tg\ ∘\ f} of the conjectured Giry-monad
lift and is reserved for follow-up work.

In dialectical-materialist language: the existence of multi-polar
contradiction structure (the full would-be \texttt{T-Alg\_T}) is not
identical with its realization in historical practice (the reachable
\texttt{T-Alg\_T\^{}\{(η,\ p\_0)\}}). The ascending phase of dialectical
motion --- the \emph{negation of negation} that produces qualitatively
new synthesis --- requires material conditions beyond those of pure
gradient descent. This provides experimental and categorical grounding
for Axiom 3's OP2.

\begin{center}\rule{0.5\linewidth}{0.5pt}\end{center}

\hypertarget{appendix-2.3.a-numerical-barrier-computation-for-a1}{%
\subsection{Appendix 2.3.A: Numerical Barrier Computation for
A1}\label{appendix-2.3.a-numerical-barrier-computation-for-a1}}

Independent SymPy verification of all critical-point quantities for
\texttt{V(u)\ =\ u(u−1)²(u−4)²}:

\begin{itemize}
\tightlist
\item
  Minima at \texttt{u\ ∈\ \{0,\ 1,\ 4\}} with \texttt{V\ =\ 0}
\item
  Inner saddle at \texttt{u\ =\ (15\ −\ √145)/10\ ≈\ 0.296},
  \texttt{V(u\_s\_in)\ =\ 2.013}
\item
  Outer saddle at \texttt{u\ =\ (15\ +\ √145)/10\ ≈\ 2.704},
  \texttt{V(u\_s\_out)\ =\ 13.187}
\item
  \texttt{V\textquotesingle{}\textquotesingle{}(u)} at minima:
  \texttt{V\textquotesingle{}\textquotesingle{}(1)\ =\ 18.00},
  \texttt{V\textquotesingle{}\textquotesingle{}(4)\ =\ 72.00}; at
  saddles:
  \texttt{V\textquotesingle{}\textquotesingle{}(0.296)\ =\ −31.41},
  \texttt{V\textquotesingle{}\textquotesingle{}(2.704)\ =\ −26.59}
\end{itemize}

Under A1's initialization
(\texttt{\textbar{}ψ\textbar{}\_init\ ≈\ 1.0\ ±\ 0.3}),
\texttt{u\_init\ ∈\ {[}0.49,\ 1.69{]}}. The inner saddle
\texttt{u\_s\_in\ =\ 0.296} is strictly below this range and the outer
saddle \texttt{u\_s\_out\ =\ 2.704} strictly above; \textbf{the entire
init support lies in the basin of attraction of \texttt{u=1}}. Under
L²-gradient flow of the free-energy functional
\texttt{F{[}ψ{]}\ =\ ∫\ V(\textbar{}ψ\textbar{}²)\ +\ \textbar{}∇ψ\textbar{}²\ dx}
(monotone \texttt{dF/dt\ ≤\ 0}; first-order dissipative; \textbf{no
kinetic-energy term}), every trajectory relaxes to \texttt{u=1} without
crossing any saddle. Observed \texttt{u\_max\ =\ 1.199} confirms.

Static energy ratio for reference:
\texttt{V(u\_s\_out)\ /\ V(u\_init\_max)\ =\ 13.187\ /\ 4.293\ ≈\ 3.07×}
(linear ratio; \textbf{not} to be confused with the dynamical
log-probability gap below).

Under overdamped Langevin (\texttt{γ\ ∂\_t\ ψ\ =\ −∂V/∂ψ*\ +\ η},
\texttt{⟨ηη*⟩\ =\ σ²δ}, \texttt{T\_eff\ =\ σ²/(2γ)}), escape becomes a
Kramers rare-event process. For \texttt{σ=0.5,\ γ=1\ →\ T\_eff=0.125},
the outer barrier gives \texttt{ΔV/T\_eff\ =\ 105.5}; Kramers escape
time \texttt{∝\ exp(105.5)} exceeds the simulation horizon
(\texttt{t\_sim=250}) by \texttt{log₁₀(exp(105.5)/250)\ ≈\ 43.4} ---
i.e.~\textbf{\textasciitilde43 orders of magnitude (logarithmic gap, not
linear ratio)}. The B1 Stage \texttt{0\%} occupancy of \texttt{u=4} is
therefore a \textbf{finite-time rare-event obstruction}, not an
equilibrium prohibition.

For the inner barrier (\texttt{ΔV=2.013,\ ΔV/T\_eff=16.1}), Kramers
prediction yields \texttt{\textasciitilde{}9.5×10⁻⁵} expected crossings
per voxel over \texttt{t\_sim=250}. Observed crossing fraction is
\texttt{0.85}, a discrepancy of factor \texttt{\textasciitilde{}10⁴}
(i.e.~\textasciitilde4 orders of magnitude). This open methodological
question is flagged in §4.8.3 and discussed in detail in the A1.4
verdict (\texttt{block4\_5/a1\_4/VERDICT.md} §5.1).

\begin{center}\rule{0.5\linewidth}{0.5pt}\end{center}

\emph{Forward references}. §3 develops the explicit PDE realization of
the adjunction \texttt{F\ ⊣\ G} for MaoField's source-field construction
and Ginzburg-Landau dynamics; Axioms 1--7 are realized as specific
structural constraints on that realization (§3.2), and the open problems
OP1 (M2 fixed-point formalization, falsified in §5.1) and OP2
(non-equilibrium extensions, refined in §5.2) are reformulated in the
concrete PDE setting. The categorical vocabulary introduced here ---
\texttt{T-algebra}, \texttt{T-Alg\_T\^{}\{(η,\ p\_0)\}},
\texttt{Δ\_OP2}, the three-faces convergence of §2.3.3 --- is referenced
throughout §§3--5 whenever the distinction between \emph{existent} and
\emph{realizable} structure matters.

\hypertarget{mathematical-formulation}{%
\section{3. Mathematical Formulation}\label{mathematical-formulation}}

\textbf{Position}: structured statement of the field-theoretic
formalism, the seven framework axioms, the source-attractor adjunction,
and the two meta-problems (OP1, OP2). Includes a dedicated subsection
(§3.5) on the use of symmetry-breaking language under finite-lattice
settings --- a concession to mathematical rigor that the empirical
sections (§4) require.

\begin{center}\rule{0.5\linewidth}{0.5pt}\end{center}

\hypertarget{field-theoretic-setup}{%
\subsection{3.1 Field-Theoretic Setup}\label{field-theoretic-setup}}

We model semantic information via a complex scalar field
\texttt{ψ\ :\ Ω\ →\ ℂ} over a discrete spatial domain
\texttt{Ω\ =\ {[}N{]}³} (default \texttt{N=32},
\texttt{\textbar{}Ω\textbar{}\ =\ 32³\ =\ 32768} voxels) with periodic
boundary conditions. The field evolves under a generalized
Ginzburg-Landau / Allen-Cahn dynamics:

\begin{verbatim}
γ · ∂_t ψ(x, t) = D · ∇²ψ(x, t) − ∂V/∂ψ*(ψ(x, t)) + S(x) + η(x, t)         (3.1)
\end{verbatim}

with parameters and quantities:

\begin{itemize}
\tightlist
\item
  \texttt{γ\ \textgreater{}\ 0}: damping coefficient (default
  \texttt{γ=1}, units fixing convention)
\item
  \texttt{D\ \textgreater{}\ 0}: diffusion coefficient (default
  \texttt{D=0.1})
\item
  \texttt{∇²}: discrete Laplacian on the toroidal lattice
  (\texttt{dx=1})
\item
  \texttt{V\ :\ ℝ\_\{≥0\}\ →\ ℝ}: scalar potential as a function of
  \texttt{u\ =\ \textbar{}ψ\textbar{}²\ =\ \textbar{}a\textbar{}²\ +\ \textbar{}b\textbar{}²}
  where \texttt{ψ\ =\ a\ +\ ib}. The Wirtinger derivative
  \texttt{∂V/∂ψ*\ =\ V\textquotesingle{}(u)\ ·\ ψ} provides the
  field-theoretic gradient (see Appendix 3.A for sign conventions and
  derivation)
\item
  \texttt{S\ :\ Ω\ →\ ℂ}: a deterministic source term encoding the input
  text. Two canonical constructions are studied:

  \begin{itemize}
  \tightlist
  \item
    \texttt{S\_sparse(x)}: byte-frequency construction (UTF-8 codepoints
    scattered across the 32³ grid; \textasciitilde10\% nonzero
    occupancy)
  \item
    \texttt{S\_dense(x)}: BGE-M3 1024-dim embedding tiled across the
    grid, smoothed and max-normalized
  \end{itemize}
\item
  \texttt{η\ :\ Ω\ ×\ ℝ\_\{≥0\}\ →\ ℂ}: complex Gaussian noise
  satisfying
  \texttt{⟨η(x,t)\ η*(x\textquotesingle{},t\textquotesingle{})⟩\ =\ σ²·δ\_\{xx\textquotesingle{}\}·δ(t−t\textquotesingle{})},
  parametrized by noise intensity \texttt{σ\ ≥\ 0}. The case
  \texttt{σ=0} recovers deterministic gradient flow
\end{itemize}

The default potential is the standard double-well GL form
\texttt{V(u)\ =\ ¼·(u\ −\ 1)²} (admitting
\texttt{\textbar{}ψ\textbar{}=1} as the unique minimum on the U(1) orbit
of vacua, and \texttt{ψ=0} as an unstable fixed point). For diagnostic
experiments (Blocks IV.5, V) we substitute richer \texttt{V(u)}; the
explicit forms are stated in §4 at point of use.

The dynamics (3.1) admits the L² free-energy functional

\begin{verbatim}
F[ψ] := ∫_Ω [ D·|∇ψ|² + V(|ψ|²) − (S·ψ* + S*·ψ) ] dx                       (3.2)
\end{verbatim}

so that \texttt{γ\ ∂\_t\ ψ\ =\ −δF/δψ*\ +\ η}. We adopt the no-½
convention in (3.2) --- equivalent to using the Wirtinger derivative
\texttt{δ/δψ*} directly without inserting the canonical Wirtinger factor
of 2 --- so that the Euler-Lagrange equation
\texttt{γ\ ∂\_t\ ψ\ =\ −δF/δψ*\ +\ η} reproduces (3.1) without
rescaling. See Appendix 3.A for the Wirtinger conventions in detail.

In the noiseless limit \texttt{σ=0},
\texttt{dF/dt\ =\ −γ⁻¹\ ∫\ \textbar{}δF/δψ*\textbar{}²\ ≤\ 0},
exhibiting \textbf{monotone descent} --- a fact we use repeatedly in §4
to diagnose what gradient flow alone cannot reach (see §4.7, §4.8).

\textbf{Numerical scheme}. Equation (3.1) is integrated by explicit
Euler with timestep \texttt{dt=0.05} for \texttt{5000} steps (simulation
horizon \texttt{t\_sim\ =\ 250}) unless noted. A clamp
\texttt{\textbar{}a\textbar{},\ \textbar{}b\textbar{}\ ≤\ 3.0} is
applied for numerical safety; in some experiments (notably §4.9 with
\texttt{σ\ ≥\ 0.5}) this clamp is reached and the resulting voxel
occupancy at the boundary is a \texttt{{[}DYNAMIC-IMPLEMENTATION{]}}
artifact rather than a physical steady state. We flag these cases
explicitly throughout §4.

\begin{center}\rule{0.5\linewidth}{0.5pt}\end{center}

\hypertarget{seven-framework-axioms}{%
\subsection{3.2 Seven Framework Axioms}\label{seven-framework-axioms}}

The MaoField framework postulates seven axioms. These are not derived;
they are postulates whose consequences (and limits) the paper explores.
Each axiom is stated, followed by its concrete realization in the
formalism of §3.1, and (where applicable) its connection to an open
meta-problem (OP1 or OP2).

\textbf{Axiom 1} --- \emph{A particle is a dynamic process, not a static
definition}.

\textbf{Realization}: The unit of representation is the field
configuration \texttt{ψ(x,\ t)} together with its evolution under (3.1),
not a static vector. The ``meaning'' of an input text is the trajectory
of \texttt{ψ} under the dynamics, evaluated at convergence (or at finite
horizon).

\textbf{Axiom 2} --- \emph{Understanding equals the equilibrium state of
internal-external opposition unified}.

\textbf{Realization}: Score-ready representation is the fixed point of
(3.1) under given source \texttt{S} and potential \texttt{V}. Adjunction
\texttt{F\ ⊣\ G} between Text and Field categories formalizes this
opposition (§2.3.1); T-algebras (objects fixed by \texttt{T\ =\ G∘F})
are the categorical form of this equilibrium (§2.3.1).

\textbf{Axiom 3} --- \emph{The driving force is the data's intrinsic
regularity, not an externally imposed objective}.

\textbf{Realization}: No external loss function or training signal. The
dynamics (3.1) is autonomous --- driven only by the source \texttt{S}
(encoding input regularity) and the potential \texttt{V} (encoding
intrinsic field structure). \textbf{Open Problem 2 (OP2)}: the fully
autonomous case (\texttt{η=0}, gradient flow only) is descending-phase
only and cannot realize the \emph{ascending} phase of dialectical motion
(§4.7, §4.8 demonstrate the empirical limits; §5.2 reformulates OP2).

\textbf{Axiom 4} --- \emph{The corpus is a record of practice, not a fit
to a statistical distribution}.

\textbf{Realization}: Source \texttt{S} is constructed deterministically
from each input text (no learned compression at the inference-time
framework boundary). \texttt{S\_sparse} (byte-level, fully transparent)
is the strict realization. \texttt{S\_dense} (BGE-M3, learned upstream
and fixed at inference time) is included as a \textbf{partial relaxation
of Axiom 4 used for ablation purposes}, not as an embodiment of the
axiom --- and the resulting source drift in §4.9
(\texttt{⟨S\_b⟩\ =\ -1.5×10⁻³,\ t\ =\ -17.1,\ p\ \textless{}\ 10⁻²⁵})
provides \textbf{quantitative evidence for the necessity of the strict
reading}: a learned embedding violates the mean-zero source assumption
underlying our Laplace prediction by 17σ, and this violation is the
dominant source of OP2's source-statistics axis (§5.2.2 (a)). The paper
studies both realizations precisely to expose this contrast.

\textbf{Axiom 5} --- \emph{Composite structures are the superposition of
simpler structures}.

\textbf{Realization}: Composite source fields
\texttt{V\_\{A∪B\}(ψ)\ =\ V\_A(ψ)\ +\ V\_B(ψ)} (potential addition, not
field addition --- GL non-linearity prevents the latter from preserving
the GL form). Sources superpose linearly:
\texttt{S\_\{A∪B\}(x)\ =\ S\_A(x)\ +\ S\_B(x)}. Compositional benchmarks
(e.g., SCAN-style) are reserved for Block V Phase B.

\textbf{Axiom 6} --- \emph{Matching is itself self-training}.

\textbf{Realization}: Given a query field \texttt{ψ\_q} and a candidate
document field \texttt{ψ\_d}, the matching score is computed by an
iteration that is jointly self-consistent --- not by a fixed encoder.
\textbf{Open Problem 1 (OP1)}: a candidate iteration
\texttt{S\_\{n+1\}\ =\ normalize(S\_n\ ·\ ψ*\_q\ ·\ ψ\_d)} (M2) was
proposed as a Banach contraction giving the unique fixed point. We
\textbf{falsify} this candidate in §5.1: M2 is 0-homogeneous and
trivially fails Banach contraction in \texttt{ℂ\^{}N}; the question of
contraction in projective \texttt{ℂP\^{}\{N-1\}} under Fubini-Study
metric remains open, but no rigorous fixed-point theorem is currently in
hand. Axiom 6 thus has \textbf{no known canonical mathematical
formalization} as of this writing.

\textbf{Axiom 7} --- \emph{Reaction rules are dynamically shaped by
input and state, not fixed}.

\textbf{Realization}: In current MaoField, \texttt{V(ψ)} is fixed across
the run; the dynamics (3.1) is the only mechanism by which input shapes
outcome. A \textbf{future direction} (Block V Phase C sub-block C-V,
dynamic-potential, see §5.4) replaces fixed \texttt{V} with
\texttt{V(ψ,\ t,\ history)}, where the potential itself is evolved by an
outer loop reflecting accumulated field history. This is consistent with
Axiom 7 but not yet implemented.

\textbf{Tension points}. We note three places where the strict reading
of an axiom comes into tension with our current implementation:

\begin{enumerate}
\def\labelenumi{\arabic{enumi}.}
\tightlist
\item
  Axiom 4 vs.~BGE source: BGE is a learned (statistically fitted)
  embedding, a partial concession against ``no statistical fit.'' We
  treat BGE as a \textbf{trained reflection of practice} --- i.e.~the
  upstream training is itself a practice-record --- but flag this as a
  concession with empirical consequences (the source drift in §4.9).
\item
  Axiom 3 vs.~Langevin noise: introducing \texttt{η} is an external
  (thermal) driving in tension with ``no externally imposed objective.''
  We treat \texttt{η} as \textbf{quasi-material} (representing
  finite-temperature physical noise rather than a designed gradient
  signal) and use it only diagnostically to expose OP2's structure
  (§4.8, §4.9). Future \emph{directed} driving (anisotropic noise,
  active forcing) would require more careful axiomatic justification.
\item
  Axiom 7 vs.~fixed \texttt{V}: as noted, full Axiom 7 implementation is
  deferred to Block V V.5.
\end{enumerate}

\begin{center}\rule{0.5\linewidth}{0.5pt}\end{center}

\hypertarget{source-attractor-adjunction}{%
\subsection{3.3 Source-Attractor
Adjunction}\label{source-attractor-adjunction}}

The dynamics (3.1) implements an adjunction \texttt{F\ ⊣\ G} between a
category of inputs and a category of attractors:

\begin{itemize}
\tightlist
\item
  \texttt{F\ :\ Text\ →\ Field} is the source construction
  \texttt{Text(s)\ ↦\ S(x)} extending to the initial field configuration
  \texttt{ψ\_0(x)}. Two implementations are studied: \texttt{F\_sparse}
  and \texttt{F\_dense} (defined in §3.1).
\item
  \texttt{G\ :\ Field\ →\ Score-ready} is the time-evolution map
  \texttt{(ψ\_0,\ V,\ σ)\ ↦\ ψ\_∞}, the fixed point (or finite-horizon
  snapshot) of (3.1).
\end{itemize}

The composition \texttt{T\_•\ :=\ G∘F\_•} becomes a well-defined
endofunctor of the ambient category \texttt{Meas} (§2.3.2). We do
\textbf{not} claim \texttt{T\_•} is a monad in the deterministic case;
the lift to a Giry-monad \texttt{P\ ∘\ T\_•} is conjectured and reserved
for follow-up work (§2.3.3 and footnote \footnote{Concretely,
  \texttt{Meas} is the category of measurable spaces with measurable
  maps; Text, Field, and Score-ready objects are embedded via their
  natural σ-algebras (discrete for Text, Borel for the ℂ\^{}N Field,
  Borel for Score). Deterministic maps \texttt{F\_•,\ G} extend to
  identity on the complement, so \texttt{T\_•\ :=\ G∘F\_•} is a bona
  fide endofunctor of \texttt{Meas}. The Giry monad \texttt{P} is the
  standard probability-measure monad on \texttt{Meas} (Giry 1982; Jacobs
  2010): for measurable space \texttt{(X,\ Σ\_X)}, \texttt{P(X)} is the
  space of probability measures on \texttt{(X,\ Σ\_X)} with the
  σ-algebra generated by evaluation maps \texttt{μ\ ↦\ μ(A)} for
  \texttt{A\ ∈\ Σ\_X}; the unit \texttt{η\_X\ :\ X\ →\ P(X)} is the
  Dirac map \texttt{x\ ↦\ δ\_x}; the multiplication
  \texttt{μ\_X\ :\ P(P(X))\ →\ P(X)} is the marginalization
  \texttt{Π\ ↦\ ∫ν\ dΠ(ν)}. The conjectured lift \texttt{P\ ∘\ T\_•} of
  our endofunctor would carry the (η, μ) of \texttt{P} provided the
  natural-transformation diagrams compose; verifying or disproving this
  is OP1's category-theoretic component, related to but distinct from
  the M2-falsification of OP1's analytic component (see §5.1).}
thereof). For the present paper, all categorical statements involving
\texttt{T\_•} are read as endofunctorial; statements requiring monad
axioms (T-algebras, Kleisli, EM) are flagged ``conjectured under monad
lift.''

A T-algebra (under the conjectured lift) \texttt{(X,\ α:\ TX\ →\ X)}
corresponds to a self-consistent attractor of the dynamics: an object
that is fixed by the iterated source-evolution-readout cycle. The
reachability question --- which T-algebras are \emph{dynamically}
attainable from a given initial distribution --- is the categorical form
of OP2 (see §2.3.4 and §5.2).

\begin{center}\rule{0.5\linewidth}{0.5pt}\end{center}

\hypertarget{open-problems-as-meta-theoretic-constraints}{%
\subsection{3.4 Open Problems as Meta-Theoretic
Constraints}\label{open-problems-as-meta-theoretic-constraints}}

Two open problems organize the framework's relationship to its own
incompleteness. Both are explicitly identified, neither is resolved
here, and both are central to the forward roadmap.

\textbf{OP1 (Axiom 6 formalization)}. Axiom 6 (``matching is
self-training'') asserts the existence of a self-consistent iteration
that produces the matching score directly from the field representations
of query and document, without an external loss. The natural candidate
iteration

\begin{verbatim}
M2:    S_{n+1} = normalize(S_n · ψ*_q · ψ_d)                                (3.3)
\end{verbatim}

was proposed as a Banach contraction giving a unique fixed point. We
falsify M2 as a contraction in §5.1. The status of OP1 is therefore:
\textbf{Axiom 6 is asserted; no known canonical mathematical
formalization exists}. The honest position is that Axiom 6 currently
functions as a \emph{physical intuition} awaiting mathematical capture,
not a theorem.

\textbf{OP2 (Axiom 3 non-equilibrium extension)}. Axiom 3 (``intrinsic
regularity drives, no external objective'') under the strict reading ---
pure gradient flow with no external driving --- is descending-phase
only: the system relaxes to the local basin of attraction of its initial
configuration and cannot realize qualitatively new synthesis (§4.7).
Closing OP2 requires extending the dynamics beyond pure gradient
descent. Section 4.9's A1.4 diagnostic refines OP2 from ``needs
non-equilibrium'' to a \textbf{three-fold co-design} problem, formalized
in §5.2.

These two open problems are \textbf{not failures of the framework but
constitutive of its research program}. A framework that announced ``all
axioms have rigorous formalization'' while implementing them in a
finite-lattice numerical experiment would be over-claiming; explicit
OP-tagging is the alternative to implicit overstatement.

\begin{center}\rule{0.5\linewidth}{0.5pt}\end{center}

\hypertarget{on-the-use-of-symmetry-breaking-language-zn-loose}{%
\subsection[3.5 On the Use of Symmetry-Breaking Language
]{\texorpdfstring{3.5 On the Use of Symmetry-Breaking Language
\footnote{See Section 3.5 (\texttt{{[}\^{}zn-loose{]}} footnote) for
  full text. Summary: \texttt{Z\_n} labels in this paper denote
  \textbf{empirical distributional remaining-symmetry patterns}, not
  rigorous group actions on a thermodynamically-broken vacuum. Strict
  SSB requires (1) group \texttt{G} acting on field space, (2)
  \texttt{V} invariant under \texttt{G}, (3) ground-state manifold is
  \texttt{G}-orbit, (4) non-trivially-transforming order parameter
  \texttt{⟨ψ⟩\ ≠\ 0}, (5) thermodynamic limit \texttt{N\ →\ ∞}.
  Conditions (3)--(5) are not verified on our \texttt{N=32³} lattice;
  \texttt{Z\_n} is a compact descriptor for convenience.}}{3.5 On the Use of Symmetry-Breaking Language }}\label{on-the-use-of-symmetry-breaking-language-zn-loose}}

A recurring formal concern arises from our repeated use of
group-theoretic labels (\texttt{U(1)}, \texttt{Z\_n}, ``spontaneous
symmetry breaking'') in describing finite-lattice simulation outputs. We
address this concern explicitly to forestall both reader confusion and
reviewer objection.

\textbf{Strict-sense spontaneous symmetry breaking} in the
Anderson-Goldstone-Nambu framework requires five conditions:

\begin{enumerate}
\def\labelenumi{\arabic{enumi}.}
\tightlist
\item
  An \textbf{order parameter} \texttt{O{[}ψ{]}} whose expectation value
  distinguishes symmetry sectors
\item
  A \textbf{symmetry group} \texttt{G} acting on field space (continuous
  or discrete)
\item
  The \textbf{Hamiltonian / action} is invariant under \texttt{G}
\item
  The \textbf{ground state} breaks \texttt{G} (existence of
  non-\texttt{G}-invariant minima)
\item
  The \textbf{thermodynamic limit}: the symmetry breaking is sharp (no
  tunneling between symmetry-related ground states) \textbf{only in the
  limit of infinite system size} (\texttt{N\ →\ ∞})
\end{enumerate}

Condition (5) is \textbf{not satisfied} by our experiments: we work at
\texttt{N\ =\ 32³\ =\ 32768} voxels, finite simulation time
\texttt{t\_sim\ =\ 250}, and finite document corpus (per-domain BEIR
scale, \texttt{O(10²\ –\ 10³)} documents). Strict-sense SSB \textbf{does
not apply}; symmetry-related sectors always have finite tunneling rate
at our scale.

\textbf{What we do mean}. We use SSB language in the \textbf{empirical
distributional sense}: a field-theoretic observable (e.g., per-document
\texttt{\textbar{}R\textbar{}} for phase concentration in §4.6,
amplitude bistability in §4.6) exhibits a distribution that, at our
finite lattice and finite simulation horizon, is \textbf{statistically
distinguishable from the symmetric (uniform / monomodal) baseline} ---
established by formal \texttt{χ²} tests in §4 with degrees of freedom
and significance levels reported.

Discrete-group labels (\texttt{Z\_n}) attached to specific regimes
(e.g., ``byte \texttt{Z\_1} regime'' in §4.6, ``BGE \texttt{Z\_2}
regime'' in §4.6) are \textbf{labels of empirical distributional
pattern}, not assertions of rigorous group action on a
thermodynamically-broken vacuum. Where we state \texttt{Z\_n} we mean:
``the empirical distribution is consistent with the residual symmetry
pattern that strict-sense SSB to \texttt{Z\_n} would produce at infinite
system size, modulo finite-size corrections we have not characterized.''

We acknowledge this distinction explicitly because conflation between
\textbf{empirical pattern} and \textbf{rigorous SSB} has historically
led to overstatement when physics formalism is applied to non-physical
(information-theoretic, biological, social) systems. \textbf{We do not
claim our finite-lattice observations have the mathematical status of
thermodynamic SSB}; we claim they are well-defined empirical
regularities under our explicit experimental protocol, suggestive of the
structural symmetry patterns we name.

\textbf{Mode-tagging convention} (used throughout §4 to reduce the
static-to-dynamic confusion that we identified internally as a recurrent
error class):

\begin{itemize}
\tightlist
\item
  \texttt{{[}STATIC{]}}: landscape geometry --- critical points,
  Hessian, barrier heights, energy ratios. No reference to time or noise
\item
  \texttt{{[}DYNAMIC-EQUILIBRIUM{]}}: Boltzmann / Laplace stationary
  distribution predictions assuming thermalization within
  \texttt{t\_sim}
\item
  \texttt{{[}DYNAMIC-RARE-EVENT{]}}: Kramers escape rate, horizon
  comparisons, \texttt{ΔV/T\_eff} regime indicators
\item
  \texttt{{[}DYNAMIC-IMPLEMENTATION{]}}: numerical scheme (CFL
  stability, integrator choice), boundary conditions, source statistics
\end{itemize}

This tagging discipline is a methodological precaution --- it makes
explicit which regime each statement belongs to and prevents the silent
conflation of, e.g., a static energy ratio with a dynamic
log-probability gap (a confusion we caught and corrected internally; see
§4.8 for the corrected statements with both quantities given
side-by-side).

\begin{quote}
\footnote{See Section 3.5 (\texttt{{[}\^{}zn-loose{]}} footnote) for
  full text. Summary: \texttt{Z\_n} labels in this paper denote
  \textbf{empirical distributional remaining-symmetry patterns}, not
  rigorous group actions on a thermodynamically-broken vacuum. Strict
  SSB requires (1) group \texttt{G} acting on field space, (2)
  \texttt{V} invariant under \texttt{G}, (3) ground-state manifold is
  \texttt{G}-orbit, (4) non-trivially-transforming order parameter
  \texttt{⟨ψ⟩\ ≠\ 0}, (5) thermodynamic limit \texttt{N\ →\ ∞}.
  Conditions (3)--(5) are not verified on our \texttt{N=32³} lattice;
  \texttt{Z\_n} is a compact descriptor for convenience.} All
\texttt{Z\_n} symbols in this paper are used in the empirical
distributional sense defined in §3.5: they describe the residual pattern
observed in the finite-lattice simulation, and are consistent with ---
but not proof of --- strict-sense Anderson-Goldstone-Nambu spontaneous
symmetry breaking. Strict-sense SSB requires the thermodynamic limit
(\texttt{N\ →\ ∞}), which is violated by our \texttt{N\ =\ 32} lattice.
This footnote applies wherever \texttt{Z\_n} appears in §4 and beyond.
\end{quote}

\begin{center}\rule{0.5\linewidth}{0.5pt}\end{center}

\hypertarget{appendix-3.a-wirtinger-derivative-convention}{%
\subsection{Appendix 3.A: Wirtinger Derivative
Convention}\label{appendix-3.a-wirtinger-derivative-convention}}

For a complex field \texttt{ψ\ =\ a\ +\ ib}, the Wirtinger derivatives
are

\begin{verbatim}
∂/∂ψ  = ½ (∂/∂a − i ∂/∂b)
∂/∂ψ* = ½ (∂/∂a + i ∂/∂b)
\end{verbatim}

For a real-valued potential \texttt{V(u)} with
\texttt{u\ =\ \textbar{}ψ\textbar{}²\ =\ a²\ +\ b²},

\begin{verbatim}
∂V/∂ψ* = V'(u) · ψ
\end{verbatim}

This is the convention used throughout the paper.

\textbf{Factor consistency between (3.1) and (3.2)}. The Wirtinger
derivative \texttt{δ/δψ*} of the no-½ free energy
\texttt{F\ =\ ∫{[}D\textbar{}∇ψ\textbar{}²\ +\ V\ −\ (Sψ*\ +\ S*ψ){]}\ dx}
(equation 3.2) is

\begin{verbatim}
δF/δψ* = −D∇²ψ + V'(u)·ψ − S
\end{verbatim}

so that
\texttt{γ\ ∂\_t\ ψ\ =\ −δF/δψ*\ =\ D∇²ψ\ −\ V\textquotesingle{}(u)·ψ\ +\ S}
reproduces (3.1) \textbf{without} any factor-of-2 rescaling. We adopt
this no-½ convention to make (3.1) and (3.2) directly compatible. The
alternative convention ---
\texttt{F̃\ =\ ∫{[}½D\textbar{}∇ψ\textbar{}²\ +\ V\ −\ ½(Sψ*\ +\ S*ψ){]}\ dx}
(with ½'s on the kinetic and source terms) --- requires the
Wirtinger-canonical relation \texttt{γ\ ∂\_t\ ψ\ =\ −2\ ·\ δF̃/δψ*} and
gives the same (3.1); we choose the no-½ convention for notational
economy. The L²-gradient flow gives, for \texttt{V(u)\ =\ ¼(u−1)²}:

\begin{verbatim}
γ ∂_t ψ = D∇²ψ − (u − 1)·ψ + S + η
\end{verbatim}

A common implementation pitfall is to use \texttt{∂V/∂ψ} (without
conjugate) in the gradient-descent direction. For radially symmetric
\texttt{V(u)}, this difference is invisible at static critical points
(where \texttt{sin} and \texttt{cos} factors of phase derivatives both
vanish) but produces incorrect dynamics in the Z\_n angular potentials
proposed for Block V (§5.3); we have verified the convention against
SymPy and against the explicit gradient of \texttt{F} defined in (3.2).

\begin{center}\rule{0.5\linewidth}{0.5pt}\end{center}

\emph{Forward references}. §4 reports the experimental evidence under
the formalism above. Block I (§4.2) establishes baseline retrieval
performance; Block II (§4.3) tests Kuramoto coupling; Block III (§4.4)
enumerates attractors; Block IV (§4.5) studies source-density × operator
interactions; Block IV.5 (§§4.6--4.9) decomposes \texttt{k*=2} into byte
and BGE regimes (Stage C), tests gradient-flow reachability of
multi-well minima (Stage A1), tests Langevin noise as an OP2 mechanism
(Stage B1), and refines OP2 to a three-fold co-design problem (Stage
A1.4). §5 then formalizes the refined open problems and lays out Block
V's experimental program.

\hypertarget{experimental-evidence}{%
\section{4. Experimental Evidence}\label{experimental-evidence}}

\textbf{Scope}: this section reports all empirical evidence for
MaoField's core claims, organized around the four canonical hypotheses
(H1--H3 + OP1/OP2 diagnostic) introduced in Section 3. All \texttt{Z\_n}
symbols are used in the \textbf{empirical distributional sense} defined
in footnote \footnote{See Section 3.5 (\texttt{{[}\^{}zn-loose{]}}
  footnote) for full text. Summary: \texttt{Z\_n} labels in this paper
  denote \textbf{empirical distributional remaining-symmetry patterns},
  not rigorous group actions on a thermodynamically-broken vacuum.
  Strict SSB requires (1) group \texttt{G} acting on field space, (2)
  \texttt{V} invariant under \texttt{G}, (3) ground-state manifold is
  \texttt{G}-orbit, (4) non-trivially-transforming order parameter
  \texttt{⟨ψ⟩\ ≠\ 0}, (5) thermodynamic limit \texttt{N\ →\ ∞}.
  Conditions (3)--(5) are not verified on our \texttt{N=32³} lattice;
  \texttt{Z\_n} is a compact descriptor for convenience.};
finite-lattice caveats (Section 3.5) apply throughout.

\textbf{Mode tags} (after Linux internal convention, following a
documented bug-prevention protocol): - \texttt{{[}STATIC{]}}: landscape
geometry (critical points, Hessian, barrier heights) -
\texttt{{[}DYNAMIC-EQUILIBRIUM{]}}: Boltzmann / Laplace stationary
distribution predictions - \texttt{{[}DYNAMIC-RARE-EVENT{]}}: Kramers
escape rate, horizon comparisons -
\texttt{{[}DYNAMIC-IMPLEMENTATION{]}}: numerical scheme, CFL, boundary
conditions, source statistics

\begin{center}\rule{0.5\linewidth}{0.5pt}\end{center}

\hypertarget{setup-and-notation}{%
\subsection{4.1 Setup and notation}\label{setup-and-notation}}

\textbf{Common configuration} across Blocks I--V.5:

\begin{itemize}
\tightlist
\item
  PDE engine: \texttt{ComplexEngine} (exp015), Ginzburg-Landau
  Allen-Cahn form
  \texttt{γ·∂\_t\ ψ\ =\ D\ ∇²ψ\ −\ ∂V/∂ψ*\ +\ S(x)\ +\ η(t)},
  \texttt{γ=1}, \texttt{D=0.1}
\item
  Default potential: \texttt{V(ψ)\ =\ ¼·(\textbar{}ψ\textbar{}²\ −\ 1)²}
  (standard GL double-well)
\item
  Grid: 3D toroidal lattice, \texttt{32³\ =\ 32768} voxels, periodic BC,
  \texttt{dx=1}
\item
  Integrator: explicit Euler, \texttt{dt=0.05} (unless noted),
  steps=5000 for reranker or 2000 for Block II
\item
  Candidate pool: BM25 top-K, K ∈ \{6, 20, 100\}
\item
  Clamp: \texttt{\textbar{}a\textbar{},\ \textbar{}b\textbar{}\ ≤\ 3.0}
  (i.e.~\texttt{\textbar{}ψ\textbar{}²\ ≤\ 18}) as numerical safeguard;
  flagged as \texttt{{[}DYNAMIC-IMPLEMENTATION{]}} artifact where it
  contributes
\item
  Datasets: BEIR subset
  \texttt{\{NFCorpus,\ SciFact,\ FiQA,\ ArguAna,\ TREC-COVID\}}; mMARCO
  and CodeSearchNet experiments exist in the exp017 archive but are not
  reported in this paper as the five-BEIR set already covers the
  hypothesis space relevant to §4
\item
  Evaluation: \texttt{nDCG@10,\ P@10,\ MRR@10}, BEIR standard
\item
  Hardware: EPYC 7B13 256-core CPU, RX 9070 XT for embeddings (bge-m3,
  bge-reranker-v2-m3 via llama.cpp)
\item
  Random seeds: 20260412 / 20260413
\end{itemize}

\textbf{Notation}: - \texttt{u\ ≡\ \textbar{}ψ\textbar{}²} (scalar
amplitude-square radial coordinate) -
\texttt{T\_eff\ ≡\ σ²/(2γ)\ =\ σ²/2} (effective temperature under
Langevin noise) - \texttt{⟨·⟩\_k} denotes ensemble mean over (documents
× voxels) in basin \emph{k}

\begin{center}\rule{0.5\linewidth}{0.5pt}\end{center}

\hypertarget{block-i-baseline-retrieval-across-5-beir-datasets}{%
\subsection{4.2 Block I --- Baseline retrieval across 5 BEIR
datasets}\label{block-i-baseline-retrieval-across-5-beir-datasets}}

We report \texttt{nDCG@10} for six scorers on five datasets to locate
MaoField's position in the retrieval performance landscape. Scorers:
BM25 (sparse), S1 (byte-ngram baseline), S4 (trigram baseline),
BGE-reranker-v2-m3 (SOTA cross-encoder), MaoField\_base (PDE with
fusion), MaoField\_E (PDE without fusion, reranker-only).

\textbf{Table 4.1}: \texttt{nDCG@10} matrix, 5 datasets × 6 scorers.

\begin{longtable}[]{@{}
  >{\raggedright\arraybackslash}p{(\columnwidth - 12\tabcolsep) * \real{0.1111}}
  >{\raggedleft\arraybackslash}p{(\columnwidth - 12\tabcolsep) * \real{0.1481}}
  >{\raggedleft\arraybackslash}p{(\columnwidth - 12\tabcolsep) * \real{0.1481}}
  >{\raggedleft\arraybackslash}p{(\columnwidth - 12\tabcolsep) * \real{0.1481}}
  >{\raggedleft\arraybackslash}p{(\columnwidth - 12\tabcolsep) * \real{0.1481}}
  >{\raggedleft\arraybackslash}p{(\columnwidth - 12\tabcolsep) * \real{0.1481}}
  >{\raggedleft\arraybackslash}p{(\columnwidth - 12\tabcolsep) * \real{0.1481}}@{}}
\toprule\noalign{}
\begin{minipage}[b]{\linewidth}\raggedright
Dataset
\end{minipage} & \begin{minipage}[b]{\linewidth}\raggedleft
BM25
\end{minipage} & \begin{minipage}[b]{\linewidth}\raggedleft
S1 byte
\end{minipage} & \begin{minipage}[b]{\linewidth}\raggedleft
S4 3-gram
\end{minipage} & \begin{minipage}[b]{\linewidth}\raggedleft
BGE-reranker
\end{minipage} & \begin{minipage}[b]{\linewidth}\raggedleft
MaoField\_base
\end{minipage} & \begin{minipage}[b]{\linewidth}\raggedleft
MaoField\_E
\end{minipage} \\
\midrule\noalign{}
\endhead
\bottomrule\noalign{}
\endlastfoot
NFCorpus & 0.3063 & 0.1366 & 0.2401 & 0.3104 & 0.0565 &
\textbf{0.2026} \\
SciFact & 0.6594 & 0.1920 & 0.4110 & 0.6223 & 0.0403 &
\textbf{0.3002} \\
FiQA & 0.2167 & 0.0686 & 0.1163 & 0.2992 & 0.0294 & \textbf{0.1082} \\
ArguAna & 0.2838 & 0.0556 & 0.2085 & 0.4560 & 0.0289 &
\textbf{0.1514} \\
TREC-COVID & 0.5589 & 0.4299 & 0.3962 & 0.7257 & 0.3553 &
\textbf{0.4932} \\
\end{longtable}

\textbf{Observations}:

\begin{enumerate}
\def\labelenumi{\arabic{enumi}.}
\tightlist
\item
  MaoField\_E vs S1 (byte baseline): 5/5 datasets positive gain,
  \textbf{+14.7\% to +172.3\%} relative improvement
\item
  MaoField\_E vs BGE-reranker: 5/5 datasets lag behind SOTA by −10.8 to
  −32.2 pp
\item
  MaoField\_E vs MaoField\_base: consistent +7.9 to +26.0 pp improvement
  from removing fusion (see §4.3 below)
\item
  BGE-reranker underperforms BM25 on SciFact by −3.7 pp, confirming SOTA
  cross-encoders have weak domains
\end{enumerate}

\textbf{Status}: MaoField\_E represents a working PDE-based reranker,
significantly above naïve byte baselines and below SOTA. This
establishes the empirical envelope within which subsequent Block
analyses operate.

\begin{center}\rule{0.5\linewidth}{0.5pt}\end{center}

\hypertarget{block-ii-kuramoto-diffusive-coupling-h1-test}{%
\subsection{4.3 Block II --- Kuramoto diffusive coupling (H1
test)}\label{block-ii-kuramoto-diffusive-coupling-h1-test}}

\textbf{Hypothesis H1} (\emph{movement to fusion}): Replacing explicit
fusion with Kuramoto-type diffusive coupling
\texttt{∂\_t\ ψ\_q\ =\ GL(ψ\_q)\ +\ g·(ψ\_d\ −\ ψ\_q)} (co-evolving
query and document fields for 2000 steps) should improve performance by
≥20\%.

\textbf{Table 4.2}: \texttt{nDCG@10} at coupling-scan optimum
\texttt{g\_best}.

\begin{longtable}[]{@{}llrrrr@{}}
\toprule\noalign{}
Dataset & Pool & \texttt{g=0} & \texttt{g\_best} & best \texttt{g} &
Relative gain \\
\midrule\noalign{}
\endhead
\bottomrule\noalign{}
\endlastfoot
NFCorpus & top-6 & 0.2201 & 0.2252 & 1.0 & +2.3\% \\
NFCorpus & top-20 & 0.1712 & 0.1712 & 0.0 & 0.0\% \\
SciFact & top-6 & 0.4082 & 0.4362 & 0.05 & +6.9\% \\
SciFact & top-20 & 0.1836 & 0.2179 & 0.05 & \textbf{+18.7\%} \\
\end{longtable}

\textbf{H1 verdict (weak falsification)}: Maximum observed gain 18.7\%,
below the pre-registered 20\% threshold. \texttt{g\_c} does not coincide
across datasets (NFC 1.0 vs SciFact 0.05). Coupling helps modestly but
does not replace fusion as hypothesized.

\textbf{Independent confirmation from Block I}: MaoField\_E (no fusion)
outperforms MaoField\_base (with fusion) by +7.9 to +26.0
pp.~\textbf{Fusion actively degrades signal} on 5/5 datasets, consistent
with the hypothesis that evolving query-document fields together erodes
discriminative information.

\begin{center}\rule{0.5\linewidth}{0.5pt}\end{center}

\hypertarget{block-iii-attractor-enumeration-h2-test}{%
\subsection{4.4 Block III --- Attractor enumeration (H2
test)}\label{block-iii-attractor-enumeration-h2-test}}

\textbf{Hypothesis H2} (\emph{contradiction, primary/secondary,
qualitative change}): The number of effective attractors \texttt{k*} is
bounded and small; this explains the top-20 vs top-100 performance
ceiling.

\textbf{Methodology}: For NFCorpus 100 docs, evolve each independently
for 5000 steps with byte source, cluster final \texttt{(a,\ b)} fields
via k-means over three representations. Select \texttt{k*} as the
silhouette-optimal cluster count over \texttt{k\ ∈\ \{2,...,5\}}.

\textbf{Table 4.3}: Silhouette-optimal attractor count.

\begin{longtable}[]{@{}lrrr@{}}
\toprule\noalign{}
Representation & Dimension & \texttt{k*} & Silhouette \\
\midrule\noalign{}
\endhead
\bottomrule\noalign{}
\endlastfoot
\texttt{raw\_ab} & 65536 & \textbf{2} & 0.214 \\
amplitude \texttt{\textbar{}ψ\textbar{}} & 32768 & \textbf{2} & 0.037 \\
phase \texttt{(cosφ,\ sinφ)} & 65536 & \textbf{2} & 0.223 \\
\end{longtable}

\textbf{H2 verdict (confirmed, with statistical caveat)}: All three
representations converge on \texttt{k*=2}. However, silhouette values
\texttt{0.037–0.223} are weak-to-moderate; \texttt{\textless{}\ 0.25} is
conventionally interpreted as ``no substantial cluster structure.'' We
thus claim \texttt{k*=2} as the \textbf{local silhouette optimum} rather
than a strongly-separated attractor count. See §4.11 for open
statistical questions.

The invariance of \texttt{k*=2} across representations motivates the
Section 4.6 diagnostic: if amplitude and phase both give \texttt{k*=2},
the underlying symmetry-breaking mechanism may yet differ.

\begin{center}\rule{0.5\linewidth}{0.5pt}\end{center}

\hypertarget{block-iv-source-density-operator-symmetry-h3-test}{%
\subsection{4.5 Block IV --- Source-density × operator-symmetry (H3
test)}\label{block-iv-source-density-operator-symmetry-h3-test}}

\textbf{Hypothesis H3} (\emph{material starting point}): Under the same
PDE operator, byte sources (materialist) and learned embedding sources
(idealist) produce different attractor structures.

\textbf{Four-quadrant setup}: byte or BGE source × PDE-dynamic or
static-cosine evaluation.

\textbf{Table 4.4}: \texttt{nDCG@10} and diagnostics across four
quadrants.

\begin{longtable}[]{@{}
  >{\raggedright\arraybackslash}p{(\columnwidth - 10\tabcolsep) * \real{0.1364}}
  >{\raggedright\arraybackslash}p{(\columnwidth - 10\tabcolsep) * \real{0.1364}}
  >{\raggedleft\arraybackslash}p{(\columnwidth - 10\tabcolsep) * \real{0.1818}}
  >{\raggedleft\arraybackslash}p{(\columnwidth - 10\tabcolsep) * \real{0.1818}}
  >{\raggedleft\arraybackslash}p{(\columnwidth - 10\tabcolsep) * \real{0.1818}}
  >{\raggedleft\arraybackslash}p{(\columnwidth - 10\tabcolsep) * \real{0.1818}}@{}}
\toprule\noalign{}
\begin{minipage}[b]{\linewidth}\raggedright
Dimension
\end{minipage} & \begin{minipage}[b]{\linewidth}\raggedright
Dataset
\end{minipage} & \begin{minipage}[b]{\linewidth}\raggedleft
IV-A byte+PDE
\end{minipage} & \begin{minipage}[b]{\linewidth}\raggedleft
IV-B BGE+PDE
\end{minipage} & \begin{minipage}[b]{\linewidth}\raggedleft
IV-C BGE cos
\end{minipage} & \begin{minipage}[b]{\linewidth}\raggedleft
IV-D byte cos
\end{minipage} \\
\midrule\noalign{}
\endhead
\bottomrule\noalign{}
\endlastfoot
Overall \texttt{nDCG@10} & NFCorpus & 0.2026 & 0.2495 & \textbf{0.3089}
& 0.2159 \\
Overall \texttt{nDCG@10} & SciFact & 0.3002 & 0.3814 & \textbf{0.6143} &
0.3555 \\
Overall \texttt{nDCG@10} & FiQA & 0.1082 & 0.1901 & \textbf{0.2774} &
0.1309 \\
Long-tail \texttt{nDCG@10} & NFC (n=227) & 0.2117 & 0.2608 &
\textbf{0.3149} & 0.2264 \\
Long-tail \texttt{nDCG@10} & FiQA (n=29) & 0.1404 & 0.1551 &
\textbf{0.2199} & 0.1341 \\
Robustness drop & SciFact ±200 & −0.010 & +0.011 & −0.015 & +0.031 \\
Attractor \texttt{k*} (raw-ab only)† & NFC 100 docs & \textbf{2} (sil
0.13) & \textbf{2} (sil 0.06) & N/A & N/A \\
\end{longtable}

† These silhouette values are for the \texttt{raw\_ab}-only clustering
in a paired IV-A vs IV-B comparison on NFCorpus 100 docs; they are
numerically lower than the Block III §4.4 triple-representation scan
(\texttt{raw\_ab} 0.214, \texttt{amplitude} 0.037, \texttt{phase} 0.223)
because Block III reports each representation independently with
per-representation optimization, while Block IV uses a common feature
space to enable cross-source comparison. Both analyses confirm
\texttt{k*=2}; only the silhouette magnitudes differ by cluster-space
choice.

‡ SciFact long-tail (n=3) row omitted from this table due to
insufficient sample; see \texttt{FINAL\_REPORT.md} §2.4 for the raw
numbers if needed.

\textbf{Observations}:

\begin{enumerate}
\def\labelenumi{\arabic{enumi}.}
\tightlist
\item
  Overall ranking consistent across three datasets:
  \texttt{IV-C\ \textgreater{}\ IV-B\ \textgreater{}\ IV-D\ \textgreater{}\ IV-A}
\item
  Within-source-type: IV-A (byte+PDE) underperforms IV-D (byte cos) by
  −1.3 to −5.5 pp; IV-B (BGE+PDE) underperforms IV-C (BGE cos) by −5.9
  to −23.3 pp
\item
  \textbf{\texttt{k*=2} holds in both IV-A and IV-B}, across
  source-density extremes
\item
  Robustness drops all \texttt{\textbar{}Δ\textbar{}\ \textless{}\ 0.04}
  --- within noise
\end{enumerate}

\textbf{H3 verdict (partial support)}: Source-density does affect
retrieval scores but does \textbf{not} produce different \texttt{k*}.
The surface-level \texttt{k*=2} identity across byte and BGE raises a
mechanistic question: does the \texttt{Z\_2} structure arise from the
\emph{same} or \emph{different} symmetry breaking in the two source
regimes? Stage C (§4.6) addresses this.

\begin{center}\rule{0.5\linewidth}{0.5pt}\end{center}

\hypertarget{stage-c-two-regimes-of-k2}{%
\subsection{\texorpdfstring{4.6 Stage C --- Two regimes of
\texttt{k*=2}}{4.6 Stage C --- Two regimes of k*=2}}\label{stage-c-two-regimes-of-k2}}

\textbf{Motivation}: \texttt{k*=2} holds for both byte and BGE sources.
If the underlying mechanism differs, Stage C's phase diagnostic will
resolve it.

\hypertarget{method}{%
\subsubsection{4.6.1 Method}\label{method}}

Compute for both IV-A and IV-B final states: - Global phase circular
mean resultant \texttt{\textbar{}R\textbar{}} (0 = uniform on S¹, 1 =
full lock) - Per-document \texttt{⟨\textbar{}R\textbar{}⟩} (average over
documents of within-document phase concentration) - 100-document
circular-mean phases, k-means on sine/cosine, measure antipodal
separation angle (expected 180° for pure \texttt{Z\_2} phase inversion)

\hypertarget{results-static}{%
\subsubsection{4.6.2 Results {[}STATIC{]}}\label{results-static}}

\textbf{Table 4.5}: Phase diagnostic for IV-A vs IV-B.

\begin{longtable}[]{@{}
  >{\raggedright\arraybackslash}p{(\columnwidth - 8\tabcolsep) * \real{0.1579}}
  >{\raggedleft\arraybackslash}p{(\columnwidth - 8\tabcolsep) * \real{0.2105}}
  >{\raggedleft\arraybackslash}p{(\columnwidth - 8\tabcolsep) * \real{0.2105}}
  >{\raggedleft\arraybackslash}p{(\columnwidth - 8\tabcolsep) * \real{0.2105}}
  >{\raggedleft\arraybackslash}p{(\columnwidth - 8\tabcolsep) * \real{0.2105}}@{}}
\toprule\noalign{}
\begin{minipage}[b]{\linewidth}\raggedright
\end{minipage} & \begin{minipage}[b]{\linewidth}\raggedleft
Global \texttt{\textbar{}R\textbar{}}
\end{minipage} & \begin{minipage}[b]{\linewidth}\raggedleft
χ²/72
\end{minipage} & \begin{minipage}[b]{\linewidth}\raggedleft
Per-doc \texttt{⟨\textbar{}R\textbar{}⟩}
\end{minipage} & \begin{minipage}[b]{\linewidth}\raggedleft
k-means centroid separation
\end{minipage} \\
\midrule\noalign{}
\endhead
\bottomrule\noalign{}
\endlastfoot
\textbf{IV-A (byte)} & \textbf{0.807} & 94989.7 & \textbf{0.840} &
\textbf{26.2°} (not antipodal) \\
\textbf{IV-B (BGE)} & 0.165 & 4879.9 & 0.197 & 81.4° (not antipodal) \\
\end{longtable}

\hypertarget{interpretation-static}{%
\subsubsection{4.6.3 Interpretation
{[}STATIC{]}}\label{interpretation-static}}

\textbf{IV-A (byte, sparse source)}:
\texttt{⟨\textbar{}R\textbar{}⟩=0.84} indicates strong phase
concentration within each document. The 100-document circular means fall
into a narrow 26.2° fan. \textbf{The U(1) phase symmetry is empirically
broken} in the sense that the observed \texttt{arg(ψ)} distribution is
statistically distinguishable from uniform (χ² ≫ 1). However, the
absence of antipodal separation (26° ≠ π rad) \textbf{rules out rigorous
\texttt{Z\_2\ ⊂\ U(1)} phase inversion} as the operative mechanism. The
pattern is better described as ``phase concentration to one preferred
direction'' --- a broken U(1) without a surviving discrete subgroup.

\textbf{IV-B (BGE, dense source)}: \texttt{⟨\textbar{}R\textbar{}⟩=0.20}
indicates much weaker phase concentration (still non-uniform; χ² = 4880
rejects uniformity strongly, but the phase modulation is weak relative
to byte). Amplitude exhibits a bimodal distribution at
\texttt{\textbar{}ψ\textbar{}∈\{0,1\}}. Note that
\texttt{V=¼(\textbar{}ψ\textbar{}²−1)²} admits only
\texttt{\textbar{}ψ\textbar{}=1} as a true minimum;
\texttt{\textbar{}ψ\textbar{}=0} is a critical point with
negative-definite Hessian (unstable fixed point, since
\texttt{V(0)=¼\ ≠\ V(1)=0}). Hence \texttt{\textbar{}ψ\textbar{}=0} and
\texttt{\textbar{}ψ\textbar{}=1} \textbf{do not lie on the same U(1)
orbit} and cannot be interchanged by any \texttt{Z\_2} group action
respecting the potential. The ``amplitude bistability'' is a bimodal
distributional pattern, not a strict \texttt{Z\_2} amplitude symmetry
breaking \footnote{See Section 3.5 (\texttt{{[}\^{}zn-loose{]}}
  footnote) for full text. Summary: \texttt{Z\_n} labels in this paper
  denote \textbf{empirical distributional remaining-symmetry patterns},
  not rigorous group actions on a thermodynamically-broken vacuum.
  Strict SSB requires (1) group \texttt{G} acting on field space, (2)
  \texttt{V} invariant under \texttt{G}, (3) ground-state manifold is
  \texttt{G}-orbit, (4) non-trivially-transforming order parameter
  \texttt{⟨ψ⟩\ ≠\ 0}, (5) thermodynamic limit \texttt{N\ →\ ∞}.
  Conditions (3)--(5) are not verified on our \texttt{N=32³} lattice;
  \texttt{Z\_n} is a compact descriptor for convenience.}.

\hypertarget{k2-in-two-regimes}{%
\subsubsection{\texorpdfstring{4.6.4 \texttt{k*=2} in two
regimes}{4.6.4 k*=2 in two regimes}}\label{k2-in-two-regimes}}

The \texttt{k*=2} identity across IV-A and IV-B hides \textbf{two
distinct distributional regimes}:

\begin{itemize}
\tightlist
\item
  Byte: phase concentration into a narrow fan; amplitude roughly
  unimodal
\item
  BGE: amplitude bimodal at \texttt{\textbar{}ψ\textbar{}∈\{0,1\}};
  phase weakly modulated
\end{itemize}

We label these ``byte Z₁ regime'' and ``BGE Z₂ regime'' in the loose
sense of footnote \footnote{See Section 3.5 (\texttt{{[}\^{}zn-loose{]}}
  footnote) for full text. Summary: \texttt{Z\_n} labels in this paper
  denote \textbf{empirical distributional remaining-symmetry patterns},
  not rigorous group actions on a thermodynamically-broken vacuum.
  Strict SSB requires (1) group \texttt{G} acting on field space, (2)
  \texttt{V} invariant under \texttt{G}, (3) ground-state manifold is
  \texttt{G}-orbit, (4) non-trivially-transforming order parameter
  \texttt{⟨ψ⟩\ ≠\ 0}, (5) thermodynamic limit \texttt{N\ →\ ∞}.
  Conditions (3)--(5) are not verified on our \texttt{N=32³} lattice;
  \texttt{Z\_n} is a compact descriptor for convenience.}, but emphasize
the labels are \textbf{distributional descriptors, not group-theoretic
claims}. The empirical finding --- that a single silhouette-optimal
\texttt{k*=2} conceals two physically distinct mechanisms --- is robust
independent of labelling convention.

\textbf{Status}: H3 receives refined support. Source density does affect
the \emph{mechanism} of attractor structure while preserving the
\emph{count}.

\begin{center}\rule{0.5\linewidth}{0.5pt}\end{center}

\hypertarget{stage-a1-multi-well-amplitude-intervention}{%
\subsection{4.7 Stage A1 --- Multi-well amplitude
intervention}\label{stage-a1-multi-well-amplitude-intervention}}

\textbf{Intervention hypothesis}: If BGE's \texttt{k*=2} arises from
amplitude bistability, introducing additional amplitude wells in
\texttt{V} should increase \texttt{k*}.

\textbf{Setup}: Replace default \texttt{V=¼(\textbar{}ψ\textbar{}²−1)²}
with triple-well \texttt{V(u)\ =\ u(u−1)²(u−4)²}, minima at
\texttt{u\ =\ \textbar{}ψ\textbar{}²\ ∈\ \{0,\ 1,\ 4\}} (i.e.,
\texttt{\textbar{}ψ\textbar{}∈\{0,\ 1,\ 2\}}). BGE-M3 embedding source,
NFCorpus 100 docs (70 with cached embeddings), init
\texttt{\textbar{}ψ\textbar{}\_init\ ≈\ 1.0\ ±\ 0.3}, all other
parameters matching IV-B.

\textbf{Stability}: 0 NaN, \texttt{u\_max=1.20}, \texttt{u\_mean=1.01},
5000 steps fully stable.

\hypertarget{amplitude-occupancy}{%
\subsubsection{4.7.1 Amplitude occupancy}\label{amplitude-occupancy}}

\textbf{Table 4.6}: Observed occupancy across three wells.

\begin{longtable}[]{@{}lllr@{}}
\toprule\noalign{}
Well & \texttt{\textbar{}ψ\textbar{}} &
\texttt{u=\textbar{}ψ\textbar{}²} & Observed occupancy \\
\midrule\noalign{}
\endhead
\bottomrule\noalign{}
\endlastfoot
Inner & 0 & 0 & \textbf{0.85\%} \\
Middle & 1 & 1 & \textbf{99.15\%} \\
Outer & 2 & 4 & \textbf{0.00\%} \\
\end{longtable}

The outer well is \textbf{not populated} under gradient flow from this
initialization.

\hypertarget{barrier-analysis-static}{%
\subsubsection{4.7.2 Barrier analysis
{[}STATIC{]}}\label{barrier-analysis-static}}

\textbf{Table 4.7}: Critical-point structure of the triple-well
potential (independent SymPy verification).

\begin{longtable}[]{@{}lrrr@{}}
\toprule\noalign{}
Point & \texttt{u} & \texttt{V} &
\texttt{V\textquotesingle{}\textquotesingle{}(u)} \\
\midrule\noalign{}
\endhead
\bottomrule\noalign{}
\endlastfoot
Min (boundary) & 0.0000 & 0.000 & --- (linear \texttt{V≈4u}) \\
Saddle (inner) & 0.2958 & \textbf{2.013} & −31.41 \\
Min & 1.0000 & 0.000 & +18.00 \\
Saddle (outer) & \textbf{2.7042} & \textbf{13.187} & −26.59 \\
Min & 4.0000 & 0.000 & +72.00 \\
\end{longtable}

\textbf{Barrier asymmetry}:
\texttt{V(saddle\ outer)\ /\ V(saddle\ inner)\ =\ 13.187\ /\ 2.013\ =\ 6.55×}.

\hypertarget{gradient-flow-inaccessibility-of-u4-dynamic-equilibrium}{%
\subsubsection{\texorpdfstring{4.7.3 Gradient-flow inaccessibility of
\texttt{u=4}
{[}DYNAMIC-EQUILIBRIUM{]}}{4.7.3 Gradient-flow inaccessibility of u=4 {[}DYNAMIC-EQUILIBRIUM{]}}}\label{gradient-flow-inaccessibility-of-u4-dynamic-equilibrium}}

Under pure gradient flow (\texttt{γ\ ·\ du/dt\ =\ −∂V/∂u}, first-order
dissipative, \textbf{no kinetic-energy term}), \texttt{V} is
monotonically non-increasing along trajectories. Initialization
\texttt{u\_init\ ∈\ {[}0.49,\ 1.69{]}} lies \textbf{entirely within the
basin of attraction of \texttt{u=1}} (inner saddle
\texttt{u\_s\_in=0.296} is left of init range, outer saddle
\texttt{u\_s\_out=2.704} is right). Hence no trajectory can reach the
\texttt{u=4} basin under gradient flow alone. Observed
\texttt{u\_max=1.199} confirms.

For reference, the \textbf{static energy ratio} is
\texttt{V(saddle\ outer)\ /\ V(u\_init\_max)\ =\ 13.187\ /\ V(1.69)\ =\ 13.187\ /\ 4.293\ ≈\ 3.07×}.
Since gradient flow is monotonically non-increasing in V and
\texttt{V(u\_init\_max)\ \textless{}\ V(saddle)}, the \texttt{u=4} basin
is \textbf{analytically unreachable} under pure gradient flow --- not an
experimental surprise but a direct consequence of the static energy
landscape geometry.

\textbf{A1 verdict}: Adding wells to \texttt{V} does not mechanically
add attractors in the observed distribution. Populating the new wells
requires either (i) initialization spanning multiple basins, (ii) a
non-equilibrium mechanism (Langevin, Hamiltonian, active), or (iii) a
structural change to \texttt{V} that brings the new minima into the
gradient-connected region. This motivates Stage B1.

\begin{center}\rule{0.5\linewidth}{0.5pt}\end{center}

\hypertarget{stage-b1-langevin-ux3c3-scan-with-kramers-timescale-analysis}{%
\subsection{4.8 Stage B1 --- Langevin σ-scan with Kramers timescale
analysis}\label{stage-b1-langevin-ux3c3-scan-with-kramers-timescale-analysis}}

\textbf{Motivation}: Test whether adding isotropic Gaussian noise
(\texttt{η(t)} with \texttt{⟨ηη*⟩\ =\ σ²δ}) enables barrier crossing to
the \texttt{u=4} basin.

\textbf{Setup}: same as A1 but with Langevin SDE
\texttt{γ\ ∂\_t\ ψ\ =\ −∂V/∂ψ*\ +\ S\ +\ η}. Scan
\texttt{σ\ ∈\ \{0,\ 0.1,\ 0.5,\ 1.0,\ 2.0\}}.

\hypertarget{observed-occupancy}{%
\subsubsection{4.8.1 Observed occupancy}\label{observed-occupancy}}

\textbf{Table 4.8}: Basin occupancy across σ-scan on triple-well
\texttt{V\ =\ u(u-1)²(u-4)²}.

\begin{longtable}[]{@{}
  >{\raggedright\arraybackslash}p{(\columnwidth - 12\tabcolsep) * \real{0.1111}}
  >{\raggedleft\arraybackslash}p{(\columnwidth - 12\tabcolsep) * \real{0.1481}}
  >{\raggedleft\arraybackslash}p{(\columnwidth - 12\tabcolsep) * \real{0.1481}}
  >{\raggedleft\arraybackslash}p{(\columnwidth - 12\tabcolsep) * \real{0.1481}}
  >{\raggedleft\arraybackslash}p{(\columnwidth - 12\tabcolsep) * \real{0.1481}}
  >{\raggedleft\arraybackslash}p{(\columnwidth - 12\tabcolsep) * \real{0.1481}}
  >{\raggedleft\arraybackslash}p{(\columnwidth - 12\tabcolsep) * \real{0.1481}}@{}}
\toprule\noalign{}
\begin{minipage}[b]{\linewidth}\raggedright
σ
\end{minipage} & \begin{minipage}[b]{\linewidth}\raggedleft
\texttt{u\_max\_all}
\end{minipage} & \begin{minipage}[b]{\linewidth}\raggedleft
\texttt{u\_mean\_all}
\end{minipage} & \begin{minipage}[b]{\linewidth}\raggedleft
near \texttt{u=0}
\end{minipage} & \begin{minipage}[b]{\linewidth}\raggedleft
near \texttt{u=1}
\end{minipage} & \begin{minipage}[b]{\linewidth}\raggedleft
near \texttt{u=4}
\end{minipage} & \begin{minipage}[b]{\linewidth}\raggedleft
clamped
\end{minipage} \\
\midrule\noalign{}
\endhead
\bottomrule\noalign{}
\endlastfoot
0.0 & 1.19 & 1.01 & 0.8\% & 99.2\% & \textbf{0.0\%} & 0.0\% \\
0.1 & 1.34 & 1.00 & 1.3\% & 98.7\% & \textbf{0.0\%} & 0.0\% \\
0.5 & 18.00 & 0.17 & \textbf{85.0\%} & 13.4\% & \textbf{0.0\%} &
0.0\% \\
1.0 & 18.00 & 17.21 & 1.9\% & 1.8\% & 0.0\% & 95.4\% \\
2.0 & 18.00 & 18.00 & 0.0\% & 0.0\% & 0.0\% & 100.0\% \\
\end{longtable}

\hypertarget{kramers-timescale-analysis-dynamic-rare-event}{%
\subsubsection{4.8.2 Kramers timescale analysis
{[}DYNAMIC-RARE-EVENT{]}}\label{kramers-timescale-analysis-dynamic-rare-event}}

\textbf{Two-regime analytical framework}: we apply the Kramers rate
formula for overdamped Langevin in 1D,

\begin{verbatim}
k_{i→j} = (1/(2π·γ))·√(V''(min_i) · |V''(saddle_ij)|) · exp(−ΔV_ij / T_eff)
\end{verbatim}

valid when \texttt{ΔV/T\_eff\ ≫\ 1} (typically ≥ 5). The \emph{expected
number of crossings} over the simulation horizon
\texttt{t\_sim\ =\ steps\ ·\ dt\ =\ 5000\ ·\ 0.05\ =\ 250} is
\texttt{k\ ×\ t\_sim}.

\textbf{Table 4.9}: Kramers analysis at σ=0.5
(\texttt{T\_eff\ =\ 0.125}), triple-well \texttt{V}.

\begin{longtable}[]{@{}
  >{\raggedright\arraybackslash}p{(\columnwidth - 10\tabcolsep) * \real{0.1364}}
  >{\raggedleft\arraybackslash}p{(\columnwidth - 10\tabcolsep) * \real{0.1818}}
  >{\raggedleft\arraybackslash}p{(\columnwidth - 10\tabcolsep) * \real{0.1818}}
  >{\raggedleft\arraybackslash}p{(\columnwidth - 10\tabcolsep) * \real{0.1818}}
  >{\raggedleft\arraybackslash}p{(\columnwidth - 10\tabcolsep) * \real{0.1818}}
  >{\raggedright\arraybackslash}p{(\columnwidth - 10\tabcolsep) * \real{0.1364}}@{}}
\toprule\noalign{}
\begin{minipage}[b]{\linewidth}\raggedright
Transition
\end{minipage} & \begin{minipage}[b]{\linewidth}\raggedleft
ΔV
\end{minipage} & \begin{minipage}[b]{\linewidth}\raggedleft
ΔV/\texttt{T\_eff}
\end{minipage} & \begin{minipage}[b]{\linewidth}\raggedleft
Prefactor
\end{minipage} & \begin{minipage}[b]{\linewidth}\raggedleft
\texttt{k·t\_sim}
\end{minipage} & \begin{minipage}[b]{\linewidth}\raggedright
Status
\end{minipage} \\
\midrule\noalign{}
\endhead
\bottomrule\noalign{}
\endlastfoot
\texttt{u=1\ →\ u=0} & 2.013 & \textbf{16.1} & \textasciitilde3.4 &
\textasciitilde8·10⁻⁵ & Kramers marginal (inner barrier) \\
\texttt{u=1\ →\ u=4} & 13.187 & \textbf{105.5} & \textasciitilde3.5 &
\textasciitilde4·10⁻⁴³ & Kramers strongly valid, \textbf{43 orders below
horizon} \\
\end{longtable}

\hypertarget{reframe-langevin-works-for-inner-fails-for-outer-7}{%
\subsubsection{\texorpdfstring{4.8.3 Reframe --- \emph{Langevin works
for inner, fails for outer}
(\#7)}{4.8.3 Reframe --- Langevin works for inner, fails for outer (\#7)}}\label{reframe-langevin-works-for-inner-fails-for-outer-7}}

The σ=0.5 observation (85\% \texttt{u=0}, 0\% \texttt{u=4}) should be
read as \textbf{two simultaneous outcomes under one noise level}:

\begin{enumerate}
\def\labelenumi{\arabic{enumi}.}
\item
  \textbf{Inner barrier} (ΔV/T\_eff = 16): Langevin-driven crossing is
  qualitatively observed. 85\% of voxels transit from the \texttt{u=1}
  basin across the inner saddle to the \texttt{u=0} basin within
  \texttt{t\_sim=250}. Note however that Kramers prediction gives only
  \textasciitilde{}\texttt{8·10⁻⁵} expected crossings per voxel --- an
  \textbf{observed/predicted ratio ≈ 1.1·10⁴ (i.e.~about 4 orders of
  magnitude)} discrepancy between 1D-Kramers and observation (see
  §4.11).
\item
  \textbf{Outer barrier} (ΔV/T\_eff = 105.5): Kramers escape time
  exceeds the simulation horizon by
  \texttt{log₁₀(exp(105.5)\ /\ 250)\ ≈\ 43.4} ---
  \textbf{\textasciitilde43 orders of magnitude} log-gap. The 0\%
  \texttt{u=4} occupancy is a \textbf{finite-time rare-event
  obstruction}, not an equilibrium prohibition.
\end{enumerate}

This reframes the original interpretation (``Langevin fails on A1'') to
a more precise dual statement:

\begin{quote}
\textbf{Langevin succeeds where the simulation horizon is adequate for
the Kramers timescale; it fails where the horizon is too short.} The
apparent ``failure'' to populate \texttt{u=4} is a timescale mismatch,
not a paradigm limitation.
\end{quote}

\textbf{Implication for OP2}: resolving OP2 (populating distant
attractors) requires either (a) \textbf{barrier-topology engineering} to
bring Kramers timescales within horizon, or (b) \textbf{directed
non-equilibrium driving} (Hamiltonian, active, anisotropic noise) to
bypass the Kramers bottleneck entirely.

\hypertarget{high-ux3c3-regime-dynamic-implementation}{%
\subsubsection{4.8.4 High-σ regime
{[}DYNAMIC-IMPLEMENTATION{]}}\label{high-ux3c3-regime-dynamic-implementation}}

At σ ≥ 1.0, the noise amplitude dominates the potential structure.
Voxels random-walk to the clamp boundary
\texttt{\textbar{}ψ\textbar{}=3} (i.e.~\texttt{u=18}), producing
\texttt{u\_max=18} and \texttt{u\_mean\ ≈\ 17} at σ=2.0. This is
\textbf{not a physical steady state} but a
\texttt{{[}DYNAMIC-IMPLEMENTATION{]}} artifact of the explicit-Euler +
clamp boundary condition combined with ΔV/T\_eff ≤ 1. See §4.11 for
methodological notes.

\begin{center}\rule{0.5\linewidth}{0.5pt}\end{center}

\hypertarget{stage-a1.4-op2-three-fold-co-design-diagnostic}{%
\subsection{4.9 Stage A1.4 --- OP2 three-fold co-design
diagnostic}\label{stage-a1.4-op2-three-fold-co-design-diagnostic}}

\textbf{Hypothesis}: If the outer barrier in A1 is too high for the
simulation horizon, lowering it should restore Laplace equilibrium under
moderate σ. Test via shifted potential
\texttt{V\_shifted(u)\ =\ u(u−1)²(u−2)²}, minima at
\texttt{u∈\{0,1,2\}}, outer barrier reduced from 13.187 (A1) to
\texttt{V(saddle\ u=1.540)\ =\ 0.095}.

\textbf{Setup}: same as A1 but with shifted \texttt{V}, σ-scan
\texttt{\{0,\ 0.1,\ 0.3,\ 0.5,\ 0.7\}} + soft-wall experiments
\texttt{V\_wall\ =\ λ·max(0,\ u−u\_max)⁴} with
\texttt{λ\ ∈\ \{1,\ 10\}}, \texttt{u\_max=2.5}, and a
\texttt{λ=1,\ σ=0.3} control.

\textbf{Laplace prediction at σ=0.5, T\_eff=0.125}
{[}DYNAMIC-EQUILIBRIUM{]}:

\begin{itemize}
\tightlist
\item
  \texttt{Z\_0\ =\ T/4\ =\ 0.0313} (linear boundary basin \texttt{V≈4u}
  near \texttt{u=0})
\item
  \texttt{Z\_1\ =\ √(2πT/V\textquotesingle{}\textquotesingle{}(1))\ =\ √(πT)\ =\ 0.6267}
  (quadratic minimum \texttt{V\textquotesingle{}\textquotesingle{}=2})
\item
  \texttt{Z\_2\ =\ √(2πT/V\textquotesingle{}\textquotesingle{}(2))\ =\ √(πT/2)\ =\ 0.4431}
  (quadratic minimum \texttt{V\textquotesingle{}\textquotesingle{}=4})
\end{itemize}

Normalization: \texttt{p\_k\ =\ Z\_k\ /\ Σ\_j\ Z\_j} with
\texttt{Σ\_j\ Z\_j\ =\ 0.0313\ +\ 0.6267\ +\ 0.4431\ =\ 1.1010}.
Predicted occupancy:
\texttt{p(0)\ ≈\ 0.028,\ p(1)\ ≈\ 0.569,\ p(2)\ ≈\ 0.402} (sum = 1.000).
Ratio \texttt{p(1)/p(2)\ =\ √2\ ≈\ 1.414} (wider basin has higher
weight). Independent rejection-sampling on \texttt{(a,b)\ ∈\ ℝ²} with
4·10⁶ points confirms this ratio.

\hypertarget{ux3c3-scan-results}{%
\subsubsection{4.9.1 σ-scan results}\label{ux3c3-scan-results}}

\textbf{Table 4.10}: Observed occupancy under σ-scan on shifted
potential (70 docs × 32³ voxels = 2.3·10⁶ samples per σ).

\begin{longtable}[]{@{}
  >{\raggedright\arraybackslash}p{(\columnwidth - 16\tabcolsep) * \real{0.0882}}
  >{\raggedleft\arraybackslash}p{(\columnwidth - 16\tabcolsep) * \real{0.1176}}
  >{\raggedleft\arraybackslash}p{(\columnwidth - 16\tabcolsep) * \real{0.1176}}
  >{\raggedleft\arraybackslash}p{(\columnwidth - 16\tabcolsep) * \real{0.1176}}
  >{\raggedleft\arraybackslash}p{(\columnwidth - 16\tabcolsep) * \real{0.1176}}
  >{\raggedleft\arraybackslash}p{(\columnwidth - 16\tabcolsep) * \real{0.1176}}
  >{\raggedleft\arraybackslash}p{(\columnwidth - 16\tabcolsep) * \real{0.1176}}
  >{\raggedleft\arraybackslash}p{(\columnwidth - 16\tabcolsep) * \real{0.1176}}
  >{\raggedright\arraybackslash}p{(\columnwidth - 16\tabcolsep) * \real{0.0882}}@{}}
\toprule\noalign{}
\begin{minipage}[b]{\linewidth}\raggedright
σ
\end{minipage} & \begin{minipage}[b]{\linewidth}\raggedleft
\texttt{T\_eff}
\end{minipage} & \begin{minipage}[b]{\linewidth}\raggedleft
\texttt{u\_max}
\end{minipage} & \begin{minipage}[b]{\linewidth}\raggedleft
\texttt{u\_mean}
\end{minipage} & \begin{minipage}[b]{\linewidth}\raggedleft
p(0)
\end{minipage} & \begin{minipage}[b]{\linewidth}\raggedleft
p(1)
\end{minipage} & \begin{minipage}[b]{\linewidth}\raggedleft
p(2)
\end{minipage} & \begin{minipage}[b]{\linewidth}\raggedleft
p(clamp)
\end{minipage} & \begin{minipage}[b]{\linewidth}\raggedright
Kramers regime
\end{minipage} \\
\midrule\noalign{}
\endhead
\bottomrule\noalign{}
\endlastfoot
0.0 & --- & 2.13 & 1.72 & 0.000 & 0.340 & \textbf{0.660} & 0.000 &
(deterministic baseline) \\
0.1 & 0.005 & 2.41 & 1.88 & 0.000 & 0.146 & \textbf{0.853} & 0.000 &
frozen (k·H≈10⁻⁷) \\
0.3 & 0.045 & 18.00 & 1.61 & 0.012 & 0.363 & 0.408 & 0.004 & partial
(k·H≈10, outer marginal) \\
0.5 & 0.125 & 18.00 & 9.90 & 0.054 & 0.246 & 0.041 & \textbf{0.525} &
ΔV/T=0.76 \textless{} 1: \textbf{Kramers invalid} \\
0.7 & 0.245 & 18.00 & 17.65 & 0.003 & 0.008 & 0.001 & \textbf{0.979} &
ΔV/T=0.39 \textless{} 1: breakdown \\
--- \textbf{Laplace} & & & & 0.028 & \textbf{0.569} & \textbf{0.402} & 0
& --- \\
\end{longtable}

\textbf{None of the σ values yield Laplace-consistent occupancy}.
Specifically:

\begin{enumerate}
\def\labelenumi{\arabic{enumi}.}
\tightlist
\item
  \textbf{σ=0 deterministic baseline} already deviates from Laplace:
  66\% \texttt{u=2} (not 40\%) even with no noise. Gradient flow plus
  source driving alone suffices to disrupt equilibrium.
\item
  \textbf{σ=0.1}: 85\% trapped in \texttt{u=2} basin. Below thermal
  mixing timescale (\texttt{k\ ·\ t\_sim\ ≈\ 10⁻⁷} for inter-well
  rates).
\item
  \textbf{σ=0.3}: ratio \texttt{p(1)/p(2)\ =\ 0.363/0.408\ =\ 0.89},
  \textbf{reversed} from Laplace prediction \texttt{1.41}.
  Detailed-balance rate ratio
  \texttt{k\_\{12\}/k\_\{21\}\ =\ 0.704\ ≈\ 1/√2} is correct, but system
  has not thermalized (21\% ``in-between''; init bias and source drift
  exceed mixing within horizon).
\item
  \textbf{σ=0.5, 0.7}: noise exceeds potential confinement; 52\% / 98\%
  of voxels hit the clamp boundary at \texttt{u=18}.
\end{enumerate}

\hypertarget{soft-wall-intervention-walled-experiments}{%
\subsubsection{4.9.2 Soft-wall intervention (walled
experiments)}\label{soft-wall-intervention-walled-experiments}}

Adding \texttt{V\_wall\ =\ λ·max(0,\ u−2.5)⁴} should confine voxels away
from the clamp.

\textbf{Table 4.11}: Walled experiments at shifted potential.

\begin{longtable}[]{@{}
  >{\raggedright\arraybackslash}p{(\columnwidth - 12\tabcolsep) * \real{0.1111}}
  >{\raggedleft\arraybackslash}p{(\columnwidth - 12\tabcolsep) * \real{0.1481}}
  >{\raggedleft\arraybackslash}p{(\columnwidth - 12\tabcolsep) * \real{0.1481}}
  >{\raggedleft\arraybackslash}p{(\columnwidth - 12\tabcolsep) * \real{0.1481}}
  >{\raggedleft\arraybackslash}p{(\columnwidth - 12\tabcolsep) * \real{0.1481}}
  >{\raggedleft\arraybackslash}p{(\columnwidth - 12\tabcolsep) * \real{0.1481}}
  >{\raggedleft\arraybackslash}p{(\columnwidth - 12\tabcolsep) * \real{0.1481}}@{}}
\toprule\noalign{}
\begin{minipage}[b]{\linewidth}\raggedright
λ
\end{minipage} & \begin{minipage}[b]{\linewidth}\raggedleft
σ
\end{minipage} & \begin{minipage}[b]{\linewidth}\raggedleft
\texttt{u\_mean}
\end{minipage} & \begin{minipage}[b]{\linewidth}\raggedleft
p(u=1)
\end{minipage} & \begin{minipage}[b]{\linewidth}\raggedleft
p(u=2)
\end{minipage} & \begin{minipage}[b]{\linewidth}\raggedleft
p(clamp)
\end{minipage} & \begin{minipage}[b]{\linewidth}\raggedleft
Δ(clamp) vs no-wall
\end{minipage} \\
\midrule\noalign{}
\endhead
\bottomrule\noalign{}
\endlastfoot
1 & 0.5 & 10.04 & 0.242 & 0.040 & 0.534 & +0.9 pp (no help) \\
10 & 0.5 & 11.04 & 0.209 & 0.031 & \textbf{0.593} & \textbf{+6.8 pp
(worse)} \\
1 & 0.3 (control) & 1.62 & 0.363 & 0.407 & 0.005 & +0.1 pp (wall
passive) \\
\end{longtable}

\textbf{Walls of this form do not restore Laplace equilibrium}. The
σ=0.3 control confirms the wall is appropriately passive when not needed
(wall active fraction 0.75\%).

\hypertarget{three-fold-diagnosis}{%
\subsubsection{4.9.3 Three-fold diagnosis}\label{three-fold-diagnosis}}

Decomposition of why the intervention fails:

\hypertarget{axis-a-source-statistics-dynamic-implementation}{%
\paragraph{Axis A --- Source statistics
{[}DYNAMIC-IMPLEMENTATION{]}}\label{axis-a-source-statistics-dynamic-implementation}}

BGE embedding's imaginary component has a \textbf{statistically extreme
non-zero mean}:

\begin{itemize}
\tightlist
\item
  Per-doc mean of 512 \texttt{Sb} components across 70 docs:
  \texttt{⟨Sb⟩\ =\ −1.48·10⁻³}, \texttt{std\ =\ 7.2·10⁻⁴}
\item
  1-sample t-test: \texttt{t\ =\ −17.12,\ p\ \textless{}\ 10⁻²⁵}
\item
  For \texttt{Sa}: \texttt{t\ =\ 0.62,\ p\ =\ 0.54} (consistent with
  zero)
\end{itemize}

A nonzero mean of \texttt{S} implies the force \texttt{−∇V\ +\ S} is
\textbf{not pure-gradient} (or at least cannot take Boltzmann form
without modification); detailed balance is broken, and the system
evolves toward a non-equilibrium steady state (NESS) rather than a
Boltzmann distribution.

Cumulative forcing scale:
\texttt{\textbar{}S\textbar{}·t\_sim\ \textasciitilde{}\ 0.1\ ·\ 250\ =\ 25},
far exceeding the shifted outer barrier \texttt{ΔV\_outer\ =\ 0.095}.
Deterministic source-pull across the outer saddle is energetically
available even at σ=0. The σ=0 baseline observation (66\% \texttt{u=2})
confirms this mechanism operates \emph{without any noise contribution}.

\hypertarget{axis-b-potential-tail-geometry-static}{%
\paragraph{Axis B --- Potential tail geometry
{[}STATIC{]}}\label{axis-b-potential-tail-geometry-static}}

Shifted \texttt{V(u)} scales as \texttt{u⁵} for large \texttt{u}
(equivalently \texttt{\textbar{}ψ\textbar{}¹⁰} in ψ-space; two
polynomial orders shallower than A1 original
\texttt{u⁹}/\texttt{\textbar{}ψ\textbar{}¹⁸}). For \texttt{σ=0.5},
per-step noise displacement is \texttt{σ√dt\ =\ 0.112} in
\texttt{(a,\ b)} coordinates, producing per-step
\texttt{Δu\ ≈\ 2\textbar{}ψ\textbar{}·0.112\ ≈\ 0.22} near basins. The
\texttt{u⁵} tail restoring force exceeds this only in a narrow band
\texttt{u\ ∈\ {[}2,\ 3{]}}; beyond, \texttt{dt·\textbar{}f\textbar{}}
overshoots and voxels random-walk to the clamp.

Soft wall \texttt{V\_wall\ =\ λ·(u−u\_max)⁴} is \texttt{C³}-smooth at
\texttt{u=u\_max} ---
\texttt{V=V\textquotesingle{}=V\textquotesingle{}\textquotesingle{}=V\textquotesingle{}\textquotesingle{}\textquotesingle{}=0},
\texttt{V\textquotesingle{}\textquotesingle{}\textquotesingle{}\textquotesingle{}=24λ}.
Wall force at \texttt{u=2.6} is only \texttt{4λ·(0.1)³\ =\ 4·10⁻³·λ}
(essentially zero); voxels transit the
\texttt{u\ ∈\ {[}u\_max,\ u\_max\ +\ δ{]}} region in
\texttt{O(δ/σ√dt)\ \textasciitilde{}\ 5} steps without appreciable
deceleration before reaching the clamp at \texttt{u=18}.

\hypertarget{axis-c-numerical-scheme-dynamic-implementation}{%
\paragraph{Axis C --- Numerical scheme
{[}DYNAMIC-IMPLEMENTATION{]}}\label{axis-c-numerical-scheme-dynamic-implementation}}

Explicit Euler stability requires
\texttt{dt·\textbar{}J\textbar{}\ \textless{}\ 2}, where \texttt{J} is
the local Jacobian spectral radius.

At \texttt{u=10,\ a²=b²=5} (isotropic case), for
\texttt{V\_wall\ =\ λ·(u−2.5)⁴}: -
\texttt{f\_wall(u)\ =\ λ\ ·\ 4\ ·\ (7.5)³\ =\ 1687.5·λ} -
\texttt{f\textquotesingle{}\_wall(u)\ =\ λ\ ·\ 12\ ·\ (7.5)²\ =\ 675·λ}
- 2×2 block
\texttt{{[}{[}J11=2a²·f\textquotesingle{}\ +\ f,\ J12=2ab·f\textquotesingle{}{]},\ {[}J12,\ J22=2b²·f\textquotesingle{}\ +\ f{]}{]}}
- At \texttt{a²=b²=5}:
\texttt{J\_max\ =\ (2a²·f\textquotesingle{}+f)\ +\ 2ab·f\textquotesingle{}\ =\ 8437.5\ +\ 6750\ =\ 15187.5·λ}

Hence: - λ=1: \texttt{J\ ≈\ 1.5·10⁴}, \texttt{dt\_max\ ≈\ 1.3·10⁻⁴},
experimental \texttt{dt=0.05} exceeds stability by
\textbf{\textsubscript{380×\textbf{ - λ=10: \texttt{J\ ≈\ 1.5·10⁵},
\texttt{dt\_max\ ≈\ 1.3·10⁻⁵}, exceeds by }}3800×}

This \textbf{directly explains} why the stronger wall (\texttt{λ=10})
\emph{worsens} clamp residency: larger Jacobian violates CFL more
severely; the clamp absorbs the instability. Explicit Euler at
\texttt{dt=0.05} is fundamentally incompatible with any meaningful wall
stiffness in this geometry.

\hypertarget{axes-are-coupled}{%
\paragraph{Axes are coupled}\label{axes-are-coupled}}

\textbf{No single-axis fix restores Laplace} (verified by elimination):

\begin{longtable}[]{@{}
  >{\raggedright\arraybackslash}p{(\columnwidth - 2\tabcolsep) * \real{0.5000}}
  >{\raggedright\arraybackslash}p{(\columnwidth - 2\tabcolsep) * \real{0.5000}}@{}}
\toprule\noalign{}
\begin{minipage}[b]{\linewidth}\raggedright
Fix axis
\end{minipage} & \begin{minipage}[b]{\linewidth}\raggedright
Remaining obstruction
\end{minipage} \\
\midrule\noalign{}
\endhead
\bottomrule\noalign{}
\endlastfoot
Source whitening only & Potential tail still \texttt{u⁵}; integrator
still CFL-violating at σ=0.5 \\
Potential shape only (deeper tail) & Source drift still breaks DB;
stiffer potential worsens integrator stability \\
Integrator only (implicit) & Source drift still \texttt{17σ} from
mean-zero; shallow tail still allows escape \\
\end{longtable}

A1.4 \textbf{refines} OP2 from a single-axis barrier-engineering problem
to a three-axis co-design problem. The three axes are empirically
orthogonal in the sense that each is an independent necessary (not
sufficient) condition for Laplace equilibrium under this class of
physics setup.

\hypertarget{a1.4-verdict}{%
\subsubsection{4.9.4 A1.4 verdict}\label{a1.4-verdict}}

A1.4 is best read as a \textbf{diagnostic refinement}, not a
success/failure binary:

\begin{itemize}
\tightlist
\item
  It \textbf{falsifies} the naive interpretation ``lowering barriers
  alone enables OP2 access''
\item
  It \textbf{does not falsify} the broader OP2 program, which still
  admits joint source/potential/integrator co-design (Block V Phase A)
\item
  It \textbf{does falsify} the specific instance (shifted potential +
  explicit Euler + BGE raw source + simple soft wall)
\item
  It provides three orthogonal engineering targets for future work
\end{itemize}

\textbf{Explicitly not claimed}: (i) OP2 is unsolvable, (ii) the
Allen-Cahn framework is broken (σ=0 baseline confirms internal dynamics
are correct), (iii) BGE embeddings are unsuitable (only their raw use as
source; whitening or gradient-projection may restore suitability).

Full details in \texttt{block4\_5/a1\_4/VERDICT.md} (companion
repository).

\begin{center}\rule{0.5\linewidth}{0.5pt}\end{center}

\hypertarget{two-regime-analytical-framework}{%
\subsection{4.10 Two-regime analytical
framework}\label{two-regime-analytical-framework}}

The preceding Block IV.5 analyses make explicit use of a
\textbf{two-regime dichotomy} for occupancy prediction:

\textbf{Regime 1 --- Rare-event (Kramers) limit}
(\texttt{ΔV/T\_eff\ ≫\ 1}): - Occupancy evolves slowly via barrier
crossings - Relevant quantity: \texttt{k·t\_sim} (expected number of
crossings) - Applicable: A1 outer barrier at σ=0.5
(\texttt{ΔV/T\ =\ 105.5}, \texttt{k·t\_sim\ ≈\ 4·10⁻⁴³} → 0 crossings)

\textbf{Regime 2 --- Thermalized (Laplace) limit}
(\texttt{ΔV/T\_eff\ ≲\ few}, system fully mixes within \texttt{t\_sim}):
- Occupancy given by equilibrium Boltzmann partition function -
\texttt{p\_k\ ∝\ Z\_k\ =\ √(2πT/V\textquotesingle{}\textquotesingle{}(u\_k))}
for quadratic minima; boundary basins require case-by-case treatment
(e.g., linear-V gives
\texttt{Z\ ∝\ T/\textbar{}V\textquotesingle{}(0+)\textbar{}}) -
Applicable: A1.4 at σ=0.5 if detailed balance held (which it does not,
per §4.9.3)

\textbf{Regime selection is dictated by the dimensionless quantity
\texttt{ΔV/T\_eff}}, not by modeling choice. Boundary regimes
(\texttt{ΔV/T\_eff\ ∈\ {[}1,\ 5{]}}) require either longer simulation or
explicit Kramers-prefactor measurement. For \texttt{ΔV/T\_eff\ ≤\ 1},
Kramers formula is invalid and system dynamics approach free diffusion
with weak potential modulation --- in the presence of boundary clamps,
this produces \texttt{{[}DYNAMIC-IMPLEMENTATION{]}} artifacts rather
than physical steady states.

Both regimes assume \texttt{⟨S⟩\ =\ 0} (mean-zero source) and standard
Boltzmann form. \textbf{§4.9.3 Axis A shows BGE violates this assumption
by \texttt{17σ}} --- a systematic caveat applicable to all MaoField
experiments in this paper that use raw BGE embeddings as source.

\begin{center}\rule{0.5\linewidth}{0.5pt}\end{center}

\hypertarget{open-methodological-questions}{%
\subsection{4.11 Open methodological
questions}\label{open-methodological-questions}}

\hypertarget{a1-inner-barrier-10ux2074-kramers-observation-discrepancy}{%
\subsubsection{4.11.1 A1 inner barrier \textasciitilde10⁴
Kramers-observation
discrepancy}\label{a1-inner-barrier-10ux2074-kramers-observation-discrepancy}}

In §4.8.3, the 1D Kramers prefactor for \texttt{u=1\ →\ u=0} gave
\texttt{k·t\_sim\ ≈\ 8·10⁻⁵} expected crossings per voxel, while
observation shows 85\% of voxels have crossed. The ratio between
prediction and observation is \textbf{factor \textasciitilde10⁴
(i.e.~\textasciitilde4 orders of magnitude)}, not 10\^{}10000.

Candidate explanations (not tested): 1. \textbf{1D Kramers prefactor
underestimates} escape rate in 3D field-theoretic setting due to
collective voxel-voxel interactions (the voxels are coupled through the
Laplacian \texttt{D∇²ψ}). To our knowledge, no closed-form
field-theoretic Kramers analog is tabulated for this geometry; empirical
calibration via σ-scan is the most direct route. 2. \textbf{Boundary
basin at \texttt{u=0}} has linear-V attractor \texttt{V≈4u} rather than
quadratic; standard Kramers prefactor derivation assumes quadratic
saddles and quadratic minima 3. \texttt{ΔV/T\_eff\ =\ 16} is marginal:
although above the conventional \texttt{≥\ 5} threshold for Kramers
validity, the \emph{absolute calibration} of the 1D prefactor in 3D
lattice field theory is not established in the literature we surveyed

\textbf{Status}: flagged as an open methodological question. The
\textbf{qualitative conclusion} (outer barrier unreachable at
\texttt{\textasciitilde{}43\ orders} log-gap) is robust regardless, but
quantitative rate matching for lower barriers requires further work ---
either empirical prefactor measurement from σ-scan, or derivation of a
field-theoretic Kramers analog.

\hypertarget{silhouette-values-0.25}{%
\subsubsection{\texorpdfstring{4.11.2 Silhouette values
\texttt{\textless{}\ 0.25}}{4.11.2 Silhouette values \textless{} 0.25}}\label{silhouette-values-0.25}}

Block III's \texttt{k*=2} finding rests on silhouette coefficients
\texttt{0.037–0.223}, below the conventional \texttt{0.25} threshold for
``substantial cluster structure''. We report \texttt{k*=2} as the
\emph{local} silhouette optimum rather than as a strongly-separated
attractor count. Future work should use alternative cluster-validation
metrics (gap statistic, bootstrap-stability) and ideally
persistent-homology diagnostics to distinguish between ``truly k-cluster
structure'' and ``k-optimum of weakly-structured point cloud''.

\hypertarget{finite-lattice-caveat}{%
\subsubsection{4.11.3 Finite-lattice
caveat}\label{finite-lattice-caveat}}

All SSB-like claims in this paper hold in the \emph{empirical
distributional sense} on a 32³ lattice within finite simulation
horizons. Strict-sense spontaneous symmetry breaking requires
thermodynamic-limit treatments (Section 3.5; footnote \footnote{See
  Section 3.5 (\texttt{{[}\^{}zn-loose{]}} footnote) for full text.
  Summary: \texttt{Z\_n} labels in this paper denote \textbf{empirical
  distributional remaining-symmetry patterns}, not rigorous group
  actions on a thermodynamically-broken vacuum. Strict SSB requires (1)
  group \texttt{G} acting on field space, (2) \texttt{V} invariant under
  \texttt{G}, (3) ground-state manifold is \texttt{G}-orbit, (4)
  non-trivially-transforming order parameter \texttt{⟨ψ⟩\ ≠\ 0}, (5)
  thermodynamic limit \texttt{N\ →\ ∞}. Conditions (3)--(5) are not
  verified on our \texttt{N=32³} lattice; \texttt{Z\_n} is a compact
  descriptor for convenience.}). The pattern separation observed at
N=32³ is suggestive of mechanism but does not constitute proof of
symmetry-sector existence in the Anderson-Goldstone-Nambu sense.

\hypertarget{missing-controls}{%
\subsubsection{4.11.4 Missing controls}\label{missing-controls}}

The following ablations were not performed and would strengthen the
empirical case: 1. σ=0 gradient-flow-only experiments for \emph{every}
σ-scan (not just A1.4) 2. Source whitening (\texttt{Sb\ −\ ⟨Sb⟩})
baseline to separate drift from mixing effects 3. Integrator-convergence
test with \texttt{dt=0.01} to rule out numerical artifacts in σ-scan
regimes 4. mMARCO cross-lingual runs (original H3 dimension, currently
SKIPPED due to download failure)

\begin{center}\rule{0.5\linewidth}{0.5pt}\end{center}

\hypertarget{experimental-status-summary}{%
\subsection{4.12 Experimental status
summary}\label{experimental-status-summary}}

\textbf{Confirmed} (empirical evidence supports):

\begin{itemize}
\tightlist
\item
  {[}H1 falsified, weak{]} Kuramoto coupling improves retrieval by max
  \texttt{+18.7\%}, below the 20\% threshold
\item
  {[}H2 confirmed, statistically weak{]} \texttt{k*=2} across three
  representations and two source types
\item
  {[}H3 partially confirmed{]} Source density affects the mechanism
  (byte vs BGE regimes) while preserving \texttt{k*}
\item
  {[}Fusion{]} Explicit query-document fusion degrades signal by
  \texttt{+7.9\ to\ +26.0\ pp} across 5 datasets
\end{itemize}

\textbf{Refined open problems}:

\begin{itemize}
\tightlist
\item
  {[}OP1{]} M2 iteration (Axiom 6 candidate) falsified in prior work;
  Axiom 6 formalization remains open
\item
  {[}OP2{]} Reframed from ``non-equilibrium extensions needed'' to
  \textbf{three-fold co-design}: source statistics (mean-zero or NESS
  treatment), potential geometry (confining tail + stability-compatible
  shape), numerical scheme (implicit or sub-CFL-dt for stiff terms).
  A1.4 diagnostic supplies all three axes; Block V Phase A is scheduled
  to test them jointly.
\end{itemize}

\textbf{Methodological} (see §4.11):

\begin{itemize}
\tightlist
\item
  A1 inner barrier Kramers prediction differs from observation by factor
  \textasciitilde10⁴
\item
  Silhouette values support \texttt{k*=2} only in local-optimum sense
\item
  BGE raw source violates mean-zero assumption of Laplace framework
\end{itemize}

\begin{center}\rule{0.5\linewidth}{0.5pt}\end{center}

\hypertarget{roadmap-refined-open-problems-and-block-v-program}{%
\section{5. Roadmap: Refined Open Problems and Block V
Program}\label{roadmap-refined-open-problems-and-block-v-program}}

\textbf{Position}: this section formalizes the two open problems (OP1,
OP2) in their post-experiment refined form, lays out Block V's
three-axis co-design experimental program (Phase A-0/A-1/A-2/A-3), and
states the three-stage publication strategy that organizes follow-up
work.

\begin{center}\rule{0.5\linewidth}{0.5pt}\end{center}

\hypertarget{op1-axiom-6-formalization-m2-falsification}{%
\subsection{5.1 OP1: Axiom 6 Formalization --- M2
Falsification}\label{op1-axiom-6-formalization-m2-falsification}}

Axiom 6 asserts that matching is self-training: the matching score
between query and document fields should be the fixed point of a
self-consistent iteration. The natural candidate

\begin{verbatim}
M2:   S_{n+1} = normalize(S_n · ψ_q* · ψ_d)                                  (reproduced from Eq. 3.3)
\end{verbatim}

was proposed as a Banach contraction giving a unique fixed point on the
unit sphere of \texttt{ℂ\^{}N}.

\textbf{Falsification}. The map
\texttt{T\_ψ\ :\ S\ ↦\ normalize(S\ ·\ ψ\_q*\ ·\ ψ\_d)} is
\textbf{0-homogeneous}: \texttt{T\_ψ(λS)\ =\ T\_ψ(S)} for any
\texttt{λ\ \textgreater{}\ 0}. Hence

\begin{verbatim}
‖ T_ψ(S) − T_ψ(2S) ‖ = ‖ T_ψ(S) − T_ψ(S) ‖ = 0
‖ S − 2S ‖           = ‖ S ‖
\end{verbatim}

The contraction coefficient is therefore identically zero between any
two points on the same ray. The Banach contraction property requires
\texttt{‖\ T(x)\ −\ T(y)\ ‖\ ≤\ k\ ·\ ‖\ x\ −\ y\ ‖} with
\texttt{k\ \textless{}\ 1} \emph{uniformly} across the metric space; M2
trivially satisfies this between rays but cannot be a contraction in any
sense that yields a unique fixed point in \texttt{ℂ\^{}N} itself (since
the entire ray is fixed setwise but no single point is fixed).

\textbf{Reformulation in \texttt{ℂP\^{}\{N-1\}} (open)}. The
0-homogeneity means \texttt{M2} descends to a well-defined map on the
projective space \texttt{ℂP\^{}\{N-1\}}. Whether the descended map is a
contraction in the \textbf{Fubini-Study metric} on
\texttt{ℂP\^{}\{N-1\}} is a meaningful question. Linear convergence
rates suggested by power-iteration analysis (Appendix A.4 of
\texttt{open\_problems\_followup.md}) are consistent with
contraction-like behavior, but a rigorous Lipschitz estimate has not
been established.

\textbf{Status}. Axiom 6 currently has \textbf{no known canonical
mathematical formalization}. M2 is excluded as a candidate. Other
candidates (M1, M3) considered in \texttt{open\_problems\_followup.md}
either fail similar tests or are demonstrably non-converging on test
cases. We \textbf{do not falsify Axiom 6 itself} --- only its M2
candidate; the axiom remains as a physical intuition awaiting
mathematical capture.

\begin{center}\rule{0.5\linewidth}{0.5pt}\end{center}

\hypertarget{op2-refined-to-three-fold-co-design}{%
\subsection{5.2 OP2: Refined to Three-Fold
Co-Design}\label{op2-refined-to-three-fold-co-design}}

\hypertarget{pre-a1.4-form}{%
\subsubsection{5.2.1 Pre-A1.4 Form}\label{pre-a1.4-form}}

Prior to A1.4 (§4.9), OP2 was stated as:

\begin{quote}
``Pure gradient flow + isotropic Langevin fails to populate distant
attractor basins under finite simulation horizon. Non-equilibrium
extensions (Langevin, Hamiltonian, active driving) are needed.''
\end{quote}

This wording treated OP2 as a single problem with potential single-axis
solution paths. The A1.4 diagnostic (§4.9) reveals this framing as
insufficient.

\hypertarget{post-a1.4-refined-form}{%
\subsubsection{5.2.2 Post-A1.4 Refined
Form}\label{post-a1.4-refined-form}}

OP2 is a \textbf{three-axis co-design problem}. Resolution in any single
axis, holding the others fixed, does not restore the predicted
equilibrium or close the reachability gap.

\textbf{OP2 (a) --- Source-field statistics}. Detailed-balance
prerequisites for Boltzmann / Laplace equilibrium require either
mean-zero sources or explicit non-equilibrium steady-state analysis.
BGE-M3 source fields exhibit
\texttt{t\ =\ -17.1,\ p\ \textless{}\ 10⁻²⁵} non-zero mean in the
imaginary component (\texttt{⟨S\_b⟩}), violating mean-zero. The
deterministic baseline \texttt{σ=0} already produces \texttt{66\%\ u=2}
occupancy (vs.~Laplace prediction \texttt{40\%}) due to source drift
alone --- \emph{no noise contribution possible}. Resolution requires
either source whitening (mean-subtraction, possibly Hodge decomposition
for solenoidal components) or an extended steady-state framework
explicitly accounting for the non-equilibrium current.

\textbf{OP2 (b) --- Potential geometry}. Confinement must exceed noise
step size across the accessible field-space domain. The shifted
potential \texttt{V(u)\ =\ u(u-1)²(u-2)²} has a \texttt{u⁵} tail at
\texttt{u\ →\ ∞}, insufficient to confine
\texttt{\textbar{}ψ\textbar{}²} against \texttt{σ\ √dt\ =\ 0.112}
per-step displacement at \texttt{σ=0.5}. Voxels escaping the outer well
random-walk to the clamp boundary (\texttt{53\%} clamp residency in
§4.9). Resolution requires either higher-order polynomial confinement
(\texttt{tail\ ≥\ u⁶}), logarithmic / exponential confinement, or ---
what we currently believe is the cleanest path ---
\textbf{non-polynomial closures} (e.g., piecewise constructions,
reflecting boundary conditions at
\texttt{\textbar{}ψ\textbar{}\ =\ R\_max}) that decouple the inter-well
barrier height from the outer confinement strength.

\textbf{OP2 (c) --- Numerical scheme}. Stiff potential gradients require
numerical integration with stability margin for the Lipschitz constant.
Soft walls \texttt{λ\ ·\ max(0,\ u\ -\ u\_max)⁴} produce Jacobian
spectral radius \texttt{J\ ≈\ 1.5×10⁴\ ·\ λ} at
\texttt{u≈u\_max\ +\ 7.5}, requiring \texttt{dt\_max,c1\ ≈\ 1.3×10⁻⁴}
(λ=1) and \texttt{dt\_max,c2\ ≈\ 1.3×10⁻⁵} (λ=10) under explicit Euler
--- \textbf{violations of \textasciitilde380× (λ=1) to
\textasciitilde3800× (λ=10)} at our \texttt{dt\ =\ 0.05} (numbers from
VERDICT.md §3.3). The result is integrator instability absorbed into the
clamp (more clamp events with stiffer wall: \texttt{λ=10} gives
\texttt{+6.8\ pp} clamp residency vs \texttt{λ=1}). Resolution requires
\textbf{implicit or IMEX integration} (wall term treated implicitly,
rest explicit), or sub-CFL \texttt{dt}, or --- what we currently believe
is the cleanest path --- \textbf{reflecting boundary conditions}
(Skorokhod radial projection) in place of soft walls.

\hypertarget{coupling-the-three-axes-are-not-independent}{%
\subsubsection{5.2.3 Coupling: The Three Axes Are Not
Independent}\label{coupling-the-three-axes-are-not-independent}}

\textbf{Fixing any axis in isolation does not restore equilibrium}:

\begin{longtable}[]{@{}
  >{\raggedright\arraybackslash}p{(\columnwidth - 2\tabcolsep) * \real{0.5000}}
  >{\raggedright\arraybackslash}p{(\columnwidth - 2\tabcolsep) * \real{0.5000}}@{}}
\toprule\noalign{}
\begin{minipage}[b]{\linewidth}\raggedright
Fix attempted
\end{minipage} & \begin{minipage}[b]{\linewidth}\raggedright
Why still fails
\end{minipage} \\
\midrule\noalign{}
\endhead
\bottomrule\noalign{}
\endlastfoot
Source whitening only & Potential tail still \texttt{u⁵}, integrator
still CFL-violating at \texttt{σ\ ≥\ 0.5} \\
Potential only (deeper tail) & Source drift still breaks detailed
balance; deeper potential worsens integrator stiffness \\
Integrator only (implicit) & Source drift remains 17σ non-zero; shallow
tail still allows escape beyond basins \\
\end{longtable}

A1.4 was an attempt to vary axis (b) alone; it produced the diagnostic
finding \textbf{precisely because} it isolated one axis and exposed the
unaddressed interactions. Block V Phase A is the systematic co-design
experimental program.

\hypertarget{connection-to-the-categorical-open-problem-2.3.3}{%
\subsubsection{5.2.4 Connection to the Categorical Open Problem
(§2.3.3)}\label{connection-to-the-categorical-open-problem-2.3.3}}

The same conjectured stochastic extension required for the Giry-monad
lift \texttt{P\ ∘\ T\_•} (§2.3.3) supplies the mathematical structure
within which OP2's three axes are naturally addressed:

\begin{itemize}
\tightlist
\item
  Source-statistics axis (OP2 a) ↔ specifying the input measure on
  \texttt{Meas} for which \texttt{T\_•} admits a Markov-kernel lift
\item
  Potential-geometry axis (OP2 b) ↔ the support and confinement
  structure of the kernel
\item
  Numerical-scheme axis (OP2 c) ↔ the discrete-time stochastic process
  approximating the continuous Markov dynamics
\end{itemize}

This identification --- three apparently independent experimental
refinements unified under a single conjectured categorical structure ---
is what we call the \textbf{three-faces convergence} of §2.3.3. Whether
the convergence holds rigorously is, of course, an open question; we
offer it as the most coherent organizing conjecture currently visible.

\begin{center}\rule{0.5\linewidth}{0.5pt}\end{center}

\hypertarget{block-v-program-phase-a-as-co-design}{%
\subsection{5.3 Block V Program: Phase A as
Co-Design}\label{block-v-program-phase-a-as-co-design}}

Block V is the next experimental campaign. Its \textbf{Phase A}
addresses OP2's three-axis structure via four sub-blocks: a reachability
pre-check (A-0) followed by axis-specific interventions (A-1, A-2, A-3)
and a joint-design evaluation (A-joint). All sub-blocks are mechanism
studies on small datasets (NFCorpus, FiQA, \textasciitilde100 docs);
larger sweeps are deferred to Phase B.

\hypertarget{a-0-reachability-pre-check}{%
\subsubsection{A-0 Reachability
Pre-Check}\label{a-0-reachability-pre-check}}

\textbf{Purpose}: Verify that the chosen
\texttt{(V,\ init,\ σ,\ t\_sim)} configuration is reachable from the
Boltzmann sense --- i.e., that the predicted equilibrium distribution is
consistent with the gradient-flow basin-of-attraction structure under
the prescribed initialization. This precludes the ``unreachable
Laplace'' trap exhibited by A1 (where the prediction was correct but the
dynamics could never reach it under any finite \texttt{σ}).

\textbf{Method}: For each candidate \texttt{V}, compute
\texttt{(BoA(V)\ ∩\ supp(p\_init))} for each minimum, verify that all
minima have nonzero overlap with init support; compute Kramers
\texttt{k\_\{i→j\}\ ·\ t\_sim} across all saddle pairs and verify they
are within a reasonable range (e.g., \texttt{{[}10⁻²,\ 10²{]}}) to
ensure thermalization is feasible within horizon.

\textbf{Output}: A pre-flight checklist for each \texttt{V}
configuration before launching the σ-scan.

\hypertarget{a-1-source-field-statistics-axis}{%
\subsubsection{A-1 Source-Field Statistics
Axis}\label{a-1-source-field-statistics-axis}}

\textbf{Purpose}: Quantify and remediate source-field bias for the
potentials studied.

\textbf{Method}: For each source construction \texttt{F\_•} (sparse
byte, dense BGE-M3, dense BGE-large, learned-but-whitened,
principled-mean-zero):

\begin{enumerate}
\def\labelenumi{\arabic{enumi}.}
\tightlist
\item
  Compute \texttt{⟨S⟩} and \texttt{t}-statistic for non-zero mean
\item
  Compute Hodge decomposition (gradient + solenoidal parts)
\item
  Test mean-subtracted variant
  \texttt{S\textquotesingle{}\ =\ S\ −\ ⟨S⟩} and re-run a baseline
  σ-scan
\item
  Compare equilibrium distributions to Laplace predictions across source
  variants
\end{enumerate}

\textbf{Success criterion}: identify a source construction (or whitening
procedure) for which the deterministic \texttt{σ=0} baseline falls
within \texttt{±10\ pp} of the gradient-flow
basin-of-attraction-weighted Laplace prediction.

\textbf{Anti-claim to avoid}: ``approaching init distribution'' (this is
not testable since init is itself a distribution).

\hypertarget{a-2-potential-geometry-axis}{%
\subsubsection{A-2 Potential Geometry
Axis}\label{a-2-potential-geometry-axis}}

\textbf{Purpose}: Find a potential family for which inter-well barrier
and outer confinement can be independently tuned.

\textbf{Sub-experiments}:

\begin{itemize}
\tightlist
\item
  \textbf{A-2.a} (deeper polynomial tail, additive form): Compare
  \texttt{V\_5(u)\ =\ u(u-1)²(u-2)²} (current A1.4) against
  \texttt{V\_8(u)\ =\ V\_5(u)\ +\ α·(u-u\_outer)\^{}k} with
  \texttt{k\ ≥\ 6} and tunable \texttt{α,\ u\_outer} adding a
  higher-order confining term --- verify whether deeper polynomial
  \emph{tail} helps without changing inter-well barrier (note:
  multiplicative tail forms like \texttt{(u-c)²} were considered and
  excluded, since they reshape inter-well barriers; see
  \texttt{BLOCK\_V\_DESIGN.md} Phase A v2 for the additive-quartic
  discussion)
\item
  \textbf{A-2.b} (Z\_n angular, full development): Implement
  \texttt{V(ψ)\ =\ -A\textbar{}ψ\textbar{}²\ +\ B\textbar{}ψ\textbar{}⁴\ +\ C·Re(ψ\^{}n)}
  for \texttt{n\ ∈\ \{3,\ 4,\ 6\}}. Confining radial structure (Mexican
  hat) decouples from angular barriers (\texttt{ε\ ·\ Re(ψ\^{}n)}). See
  \texttt{BLOCK\_V\_DESIGN.md} §2.3 for full mathematical analysis
  (Linux Claude, 2026-04-14)
\end{itemize}

(Note: the piecewise / non-polynomial route originally considered as
A-2.c has been merged into A-3.c --- Skorokhod radial reflection ---
since reflecting BC is the cleanest non-PDE alternative and properly
belongs to the numerical-scheme axis.)

\textbf{Success criterion}: identify a \texttt{V} family for which
Laplace prediction matches σ=0.3 / σ=0.5 observation within
\texttt{±15\ pp} per basin \textbf{after} A-1 source whitening.

\hypertarget{a-3-numerical-scheme-axis}{%
\subsubsection{A-3 Numerical Scheme
Axis}\label{a-3-numerical-scheme-axis}}

\textbf{Purpose}: Implement integration scheme appropriate for stiff
potentials.

\textbf{Sub-experiments}:

\begin{itemize}
\tightlist
\item
  \textbf{A-3.a} (IMEX): Treat wall term implicitly (Crank-Nicolson or
  backward Euler on the wall part), rest explicitly. Verify stability at
  \texttt{dt\ =\ 0.05} for \texttt{λ\ ∈\ \{1,\ 10,\ 100\}}
\item
  \textbf{A-3.b} (specify Jacobian structure): For implicit treatment,
  the Jacobian of \texttt{−V\textquotesingle{}(u)·ψ} is a 2×2 block per
  voxel; structure documented in \texttt{BLOCK\_V\_DESIGN.md} §4.5.3
  Phase A-3 (Linux, 2026-04-14)
\item
  \textbf{A-3.c} (Skorokhod radial reflection): Replace soft wall with a
  reflecting BC at \texttt{\textbar{}ψ\textbar{}\ =\ R\_max} via
  Skorokhod radial projection (overdamped systems have no velocity to
  reflect; the projection is done on position). This is the structurally
  cleanest BC for confined diffusion
\end{itemize}

\textbf{Success criterion}: stable integration at \texttt{dt\ =\ 0.05}
for arbitrary \texttt{V} configurations encountered in A-2.

\hypertarget{a-joint-phase-a-capstone}{%
\subsubsection{A-joint Phase A
Capstone}\label{a-joint-phase-a-capstone}}

\textbf{Purpose}: Demonstrate that the three axes addressed jointly
recover a Laplace-predictable equilibrium.

\textbf{Method}: Take the \textbf{best A-1 source configuration}
(mean-subtracted BGE or principled mean-zero), the \textbf{best A-2
potential} (selected from A-2 sub-experiments by success criterion), and
the \textbf{best A-3 integrator} (validated stable in A-3). Run σ-scan
on this joint configuration, compare to Laplace prediction.

\textbf{Success criterion}: at least one σ value yields all-basin
occupancy within \texttt{±10\ pp} of Laplace prediction \emph{with}
\texttt{{[}STATIC{]}} and \texttt{{[}DYNAMIC-EQUILIBRIUM{]}} mode tags
consistent.

\textbf{Anti-claim to avoid}: ``best of'' combinatorial over all (A-1,
A-2, A-3) variants would explode; A-joint commits to one variant per
axis as decided by the per-axis success criteria.

\hypertarget{a-joint-outcomes-and-their-interpretation}{%
\subsubsection{A-joint outcomes and their
interpretation}\label{a-joint-outcomes-and-their-interpretation}}

\begin{itemize}
\tightlist
\item
  \textbf{A-joint succeeds} → OP2 is \textbf{confirmed solvable in the
  gradient-flow + isotropic-Langevin regime} with three-axis co-design.
  OP2 (b) directed-driving program becomes a generalization rather than
  a necessity. Block V can proceed to Phase B compositional benchmarks
\item
  \textbf{A-joint partially succeeds} (e.g., 2 of 3 basins within
  tolerance) → OP2 has a \textbf{partial co-design solution}; the
  unresolved third axis identifies the remaining open question, which
  may require Phase B's directed-driving (Hamiltonian flow, anisotropic
  noise) approaches
\item
  \textbf{A-joint fails} → strong evidence that gradient-flow +
  isotropic-Langevin is \textbf{structurally insufficient} for our class
  of problems; OP2 (b) becomes the only path forward; Block V Phase B
  becomes the central thrust
\end{itemize}

In all three cases, the outcome is informative --- A-joint is a
diagnostic, not a make-or-break engineering test.

\begin{center}\rule{0.5\linewidth}{0.5pt}\end{center}

\hypertarget{block-v-phase-b-and-c-outlook}{%
\subsection{5.4 Block V Phase B and C:
Outlook}\label{block-v-phase-b-and-c-outlook}}

\textbf{Phase B} addresses OP2 path (b) --- directed non-equilibrium
driving:

\begin{itemize}
\tightlist
\item
  \textbf{B-H (Hamiltonian)}: Add symplectic structure
  \texttt{∂\_t\ ψ\ =\ -i\ δH/δψ*} as an alternative to \texttt{-δH/δψ*}.
  Energy-conserving dynamics may explore the energy shell rather than
  descend
\item
  \textbf{B-A (Active forcing)}: External drive \texttt{f(ψ,\ t)}
  representing ``practical intervention'' beyond thermal noise
\item
  \textbf{B-N (Anisotropic noise)}: \texttt{⟨η\ η*⟩} with non-trivial
  spatial / spectral structure to bypass timescale problems via directed
  exploration
\end{itemize}

\textbf{Phase C} addresses Axiom 7 (dynamic potential) and
compositionality:

\begin{itemize}
\tightlist
\item
  \textbf{C-V} (dynamic V): \texttt{V(ψ,\ t,\ history)}, the potential
  itself evolved by an outer loop reflecting field history
\item
  \textbf{C-S} (SCAN-like compositional benchmarks): test whether
  MaoField's symbolic-composition properties scale to compositional
  generalization tests once Phase A or B has stabilized the underlying
  dynamics
\end{itemize}

Phases B and C are scoped at \textbf{2026-Q3 / Q4} in our internal
planning; arXiv v1 documents Phase A only, with B and C signposted as
future work.

\begin{center}\rule{0.5\linewidth}{0.5pt}\end{center}

\hypertarget{three-stage-publication-plan}{%
\subsection{5.5 Three-Stage Publication
Plan}\label{three-stage-publication-plan}}

Internal scheduling, included for transparency:

\textbf{Stage 1 --- arXiv v1 / Zenodo v0.1.0 (2026-04-13 / 2026-04-20)}

\begin{itemize}
\tightlist
\item
  v0.1.0 skeleton release on Zenodo (DOI
  \texttt{10.5281/zenodo.19550342}): repository structure, axioms,
  narrative, no paper PDF (already released)
\item
  v0.1.1 paper PDF release: this paper (target 2026-04-20). Includes
  Block I-IV.5 experimental evidence, OP1 / OP2 statements, Phase A
  program
\end{itemize}

\textbf{Stage 2 --- Mechanism Paper at IPM (target 2026-07--08 month)}

\begin{itemize}
\tightlist
\item
  Single-purpose paper: \emph{``Why PDE-based retrieval scales poorly:
  an attractor-capacity analysis''}
\item
  Focus: mechanism explanation of \texttt{k*\ =\ 2} ceiling and
  source-density coupling
\item
  Audience: Information Processing \& Management readership (IR /
  IR-adjacent ML)
\item
  Builds on this arXiv v1's Block I-IV.5 results + Phase A outcomes
  (pending)
\end{itemize}

\textbf{Stage 3 --- Capstone Paper at Generalist Top-Tier (target 2027
Q1-Q2)}

\begin{itemize}
\tightlist
\item
  Audience: \texttt{Nature\ Machine\ Intelligence} or equivalent
\item
  Content: full framework synthesis, philosophical positioning, capstone
  results from Block V Phases A + B (and, if successful, C)
\item
  Deliberately deferred to allow Block V experimental program to mature;
  deferral is a \textbf{methodological decision}, not a timeline slip
\end{itemize}

\textbf{Why three stages}. Releasing a generalist top-tier paper before
the mechanism is empirically validated would either over-claim or
under-claim. The arXiv v1 establishes the framework and open problems
publicly; the IPM paper validates the central mechanism in a
domain-specific venue; the generalist paper integrates after validation.
This sequencing also reduces single-point publication risk (failure or
delay of any one stage does not block the others).

\begin{center}\rule{0.5\linewidth}{0.5pt}\end{center}

\hypertarget{concluding-note-on-the-status-of-op1-and-op2}{%
\subsection{5.6 Concluding Note on the Status of OP1 and
OP2}\label{concluding-note-on-the-status-of-op1-and-op2}}

We close with an explicit statement of the two open problems' present
status, intended to forestall over-reading of the contributions claimed
in this paper.

\begin{itemize}
\tightlist
\item
  \textbf{OP1 (Axiom 6 formalization)}: candidate iteration M2 falsified
  (§5.1). Axiom 6 remains as a physical intuition without canonical
  mathematical formalization. We do not claim OP1 is solved or refuted;
  we claim one candidate is excluded.
\item
  \textbf{OP2 (Axiom 3 non-equilibrium extension)}: refined from ``needs
  non-equilibrium'' to a three-axis co-design problem (§5.2). Block V
  Phase A is the experimental program designed to address it (§5.3). We
  do not claim OP2 is solved; we claim its scope is more precisely
  identified than at the start of this campaign.
\end{itemize}

\textbf{A framework whose research program lies entirely in its open
problems is not an evasion} --- it is the honest representation of a
research stage. We expose OP1 and OP2 prominently rather than concealed
because (i) the contributions of §4 (empirical evidence) and §2
(categorical structure) make sense only against a scaffold that includes
these problems explicitly, and (ii) the dialectical-materialist
epistemology that motivates the framework demands that ``open
contradictions are the engine of motion'' --- a methodological
self-application of Axiom 1.

The next paper in the series (the IPM mechanism paper, Stage 2) will
report Block V Phase A's outcome under the success criteria stated in
§5.3.

\hypertarget{discussion}{%
\section{6. Discussion}\label{discussion}}

\begin{center}\rule{0.5\linewidth}{0.5pt}\end{center}

\hypertarget{summary-of-what-this-paper-claims-and-does-not-claim}{%
\subsection{6.1 Summary of What This Paper Claims and Does Not
Claim}\label{summary-of-what-this-paper-claims-and-does-not-claim}}

MaoField is a research-stage framework, not a completed theory. We
summarize explicitly.

\textbf{Claimed}: - A categorical structure (Lawvere endofunctor on
\texttt{Meas}, conjectured Giry-monad lift) that formalizes dialectical
motion in PDE-based AI (§2.3) - A source-density decomposition that
explains the \texttt{k*\ =\ 2} retrieval ceiling via two distinct
distributional regimes (§4.6) - A refined formulation of OP2 (Axiom 3's
non-equilibrium extension) from a single-axis ``directed driving''
problem to a three-axis co-design problem with concrete sub-problem
statements (§5.2) - A working PDE-based retrieval reranker that exceeds
byte-level lexical baselines by \texttt{+14.7\%\ to\ +172.3\%} across
five BEIR datasets while lagging SOTA cross-encoders by
\texttt{−10.8\ to\ −32.2\ pp} (§4.2) - An exclusion result for Axiom 6:
the principal candidate iteration (M2) of Axiom 6's mathematical
formalization is falsified as a Banach contraction in \texttt{ℂ\^{}N} by
0-homogeneity (§5.1). The axiom itself remains open (OP1);
methodological pragmatics (mode-tagging discipline, spawn-agent review
pipeline, \footnote{See Section 3.5 (\texttt{{[}\^{}zn-loose{]}}
  footnote) for full text. Summary: \texttt{Z\_n} labels in this paper
  denote \textbf{empirical distributional remaining-symmetry patterns},
  not rigorous group actions on a thermodynamically-broken vacuum.
  Strict SSB requires (1) group \texttt{G} acting on field space, (2)
  \texttt{V} invariant under \texttt{G}, (3) ground-state manifold is
  \texttt{G}-orbit, (4) non-trivially-transforming order parameter
  \texttt{⟨ψ⟩\ ≠\ 0}, (5) thermodynamic limit \texttt{N\ →\ ∞}.
  Conditions (3)--(5) are not verified on our \texttt{N=32³} lattice;
  \texttt{Z\_n} is a compact descriptor for convenience.} convention)
are discussed separately in §6.3 as working commitments, not as
paper-level claims

\textbf{Not claimed}: - That MaoField replaces or competes with
deep-learning-based retrieval on SOTA benchmarks - That dialectical
materialism is \emph{uniquely} realized by our categorical formalism;
other formalisms may admit dialectical readings - That OP1 or OP2 is
solved; we report one candidate excluded (M2 in OP1) and one problem
scope refined (OP2 three-fold), nothing stronger - That the Giry-monad
lift conjectured in §2.3.3 is proven; it is the most coherent organizing
conjecture currently visible, not a theorem - That our finite-lattice
(\texttt{N\ =\ 32³}) observations have the status of thermodynamic-limit
spontaneous symmetry breaking; they are well-defined empirical patterns
consistent with --- but not proof of --- strict-sense SSB (§3.5)

The distinction between these two lists is, in our view, the operative
integrity check for a research-stage paper. Papers that conflate them
--- claiming the former while implicitly asserting the latter --- are
the recurring failure mode in over-ambitious framework work. We have
tried to keep the two lists separate.

\hypertarget{limitations}{%
\subsection{6.2 Limitations}\label{limitations}}

The framework's current limitations fall into four categories.

\hypertarget{empirical-limitations}{%
\subsubsection{6.2.1 Empirical
Limitations}\label{empirical-limitations}}

Our retrieval evaluation is limited to five BEIR datasets at modest
scale (per-domain \texttt{O(10²\ –\ 10³)} documents). We have not
tested: - Compositional generalization (SCAN-style benchmarks) -
Long-context coherence (beyond single-document scoring) - Cross-lingual
performance - Adversarial / out-of-distribution robustness

Block V (§5.3--5.4) is designed in part to address these, but the
results are not in this paper. A reader seeking a full empirical
envelope should read this paper as establishing the \textbf{mechanism}
and awaiting the follow-up IPM paper (§5.5 Stage 2) for broader
performance validation.

\hypertarget{scale-limitations}{%
\subsubsection{6.2.2 Scale Limitations}\label{scale-limitations}}

All field-theoretic simulations run on \texttt{32³\ =\ 32768} voxels per
document field. Strict-sense symmetry breaking requires \texttt{N\ →\ ∞}
(§3.5); at our scale, symmetry-related sectors have finite tunneling
rates that we have not fully characterized. Whether the observed
phenomena (notably the \texttt{k*\ =\ 2} mechanism decomposition in
§4.6) are stable under scale increase is an empirical question that this
paper does not resolve. Block V Phase A-0 (§5.3) includes a reachability
pre-check that partially addresses scale-related artifacts; a systematic
finite-size scaling study is reserved for follow-up work.

\hypertarget{methodological-limitations}{%
\subsubsection{6.2.3 Methodological
Limitations}\label{methodological-limitations}}

The A1 inner barrier Kramers discrepancy (observed/predicted ratio
\textasciitilde10⁴, §4.8.3 and A1.4 VERDICT §5.1) indicates that our 1D
Kramers rate predictions may systematically underestimate escape rates
in 3D lattice field theory. We flag this as an open methodological
question; our qualitative conclusions (outer barrier at
43-orders-of-magnitude log-gap) are robust, but quantitative rate
matching for near-Kramers-regime barriers requires a field-theoretic
analog to the 1D Kramers formula. This is potentially addressable via
empirical prefactor calibration from the σ-scan, deferred to Block V
Phase A-1.

\hypertarget{formal-limitations}{%
\subsubsection{6.2.4 Formal Limitations}\label{formal-limitations}}

The Lawvere-monad framing in §2.3 is \textbf{position-paper-level, not
formalization-level}. Three gaps are explicit:

\begin{enumerate}
\def\labelenumi{\arabic{enumi}.}
\tightlist
\item
  The endofunctor \texttt{T\_•} on \texttt{Meas} is rigorous, but the
  Giry-monad lift \texttt{P\ ∘\ T\_•} is conjectured (§2.3.3 footnote
  \footnote{Concretely, \texttt{Meas} is the category of measurable
    spaces with measurable maps; Text, Field, and Score-ready objects
    are embedded via their natural σ-algebras (discrete for Text, Borel
    for the ℂ\^{}N Field, Borel for Score). Deterministic maps
    \texttt{F\_•,\ G} extend to identity on the complement, so
    \texttt{T\_•\ :=\ G∘F\_•} is a bona fide endofunctor of
    \texttt{Meas}. The Giry monad \texttt{P} is the standard
    probability-measure monad on \texttt{Meas} (Giry 1982; Jacobs 2010):
    for measurable space \texttt{(X,\ Σ\_X)}, \texttt{P(X)} is the space
    of probability measures on \texttt{(X,\ Σ\_X)} with the σ-algebra
    generated by evaluation maps \texttt{μ\ ↦\ μ(A)} for
    \texttt{A\ ∈\ Σ\_X}; the unit \texttt{η\_X\ :\ X\ →\ P(X)} is the
    Dirac map \texttt{x\ ↦\ δ\_x}; the multiplication
    \texttt{μ\_X\ :\ P(P(X))\ →\ P(X)} is the marginalization
    \texttt{Π\ ↦\ ∫ν\ dΠ(ν)}. The conjectured lift \texttt{P\ ∘\ T\_•}
    of our endofunctor would carry the (η, μ) of \texttt{P} provided the
    natural-transformation diagrams compose; verifying or disproving
    this is OP1's category-theoretic component, related to but distinct
    from the M2-falsification of OP1's analytic component (see §5.1).})
\item
  The reachable-subcategory definition \texttt{T-Alg\^{}\{(η,\ p\_0)\}}
  uses a ``Kleisli-reachable path'' notion that requires the conjectured
  monad lift for full rigor (§2.3.4)
\item
  \texttt{Δ\_OP2\ :=\ \textbar{}T-Alg\textbar{}\ ⊖\ \textbar{}T-Alg\^{}\{(η,\ p\_0)\}\textbar{}}
  uses cardinality complement as a symbolic first pass; alternative
  quantifications (categorical entropy, persistent-homology dimension,
  Kan extension obstructions) may disagree in finer regimes
\end{enumerate}

A follow-up paper dedicated to the categorical formalization --- one in
which (1), (2), (3) are either resolved or their independence is
precisely characterized --- is the natural sequel to the present work.

\hypertarget{methodological-reflection}{%
\subsection{6.3 Methodological
Reflection}\label{methodological-reflection}}

This paper was produced via a workflow that deserves explicit
methodological reflection, because the workflow itself embodies
commitments that the framework argues for.

\textbf{Spawn-agent review pipeline.} During the 2026-04-13--14 campaign
in which this paper was drafted, the author working with two AI
assistants instituted a standing rule: every significant deliverable
(mathematical derivation, experimental verdict, paper section) is
submitted to an \textbf{independent paper-review subagent} before being
finalized. The reviewer is a separate instance without the drafting
agent's context, prompted with explicit review criteria (P0 must-fix, P1
recommended, P2 suggested weakening, P3 already-honest).

Under this protocol, \textbf{11 substantive errors} were caught and
corrected during Phase 1 (Kramers formulation errors,
Wirtinger-derivative bugs, over-claimed symmetry-breaking labels,
Kramers vs.~Boltzmann regime conflation, and others). A further
\textbf{pass of reviews} on the arXiv v1 section drafts (the present
paper) caught additional issues (factor-of-2 inconsistency between field
equation and free energy, Axiom 4 softening that undercut quantitative
evidence, cross-reference naming drift, CFL violation number
inconsistency). Each of these would have been noticed by a careful human
reviewer; the protocol makes the noticing \textbf{systematic rather than
contingent on reviewer attention}.

We offer this as a concrete practice compatible with the framework's
epistemology: \textbf{the dialectical-materialist insistence on
practice-driven cognition maps directly onto a ``review-as-practice''
discipline for paper production}. The reviewer is not a decorative
safety net; it is the practice that tests the drafting agent's claims
against independent analysis.

\textbf{Mode-tagging discipline.} Throughout the experimental sections
we tag statements as \texttt{{[}STATIC{]}} (landscape geometry),
\texttt{{[}DYNAMIC-EQUILIBRIUM{]}} (Boltzmann / Laplace predictions),
\texttt{{[}DYNAMIC-RARE-EVENT{]}} (Kramers escape), or
\texttt{{[}DYNAMIC-IMPLEMENTATION{]}} (integrator, boundary, source
statistics). This discipline was introduced in response to a recurring
error class identified during Phase 1 review: \textbf{static analysis
tools (critical-point geometry, Hessian analysis) being silently
transported to dynamic prediction tasks (occupancy, escape rates)
without changing toolkit}. Tagging makes the regime explicit, forcing
the author to consciously choose the appropriate formula (Kramers rate
vs.~Laplace equilibrium vs.~finite-time simulation) rather than reach
for the most convenient.

\textbf{Explicit OP-signposting.} Axioms 3 and 6 are maintained as open
problems (OP2 and OP1 respectively) rather than being implicitly
asserted or silently weakened. This is, again, a dialectical-materialist
methodological commitment: open contradictions should be made visible
rather than papered over, because their resolution is the engine of the
research program's motion.

\hypertarget{implications-for-the-ai-paradigm-discussion}{%
\subsection{6.4 Implications for the AI Paradigm
Discussion}\label{implications-for-the-ai-paradigm-discussion}}

We do not claim that MaoField refutes, replaces, or dominates
statistical deep learning. We claim that a \textbf{second paradigm} ---
distinct in motivation, mathematical infrastructure, and implementation
details --- is viable at a research-stage level and produces properties
that the first paradigm has not produced (end-to-end structural
interpretability, autonomous dynamics, explicit attractor enumeration).
The paradigm-level scope is deliberate: retrieval is the
\textbf{tractable test bed} on which §4 demonstrates non-vacuity, but
the framework's target domain is any setting where autonomous structural
dynamics rather than statistical prediction is the operative goal.

\textbf{The paradigms are complementary, not competitive}. Statistical
deep learning is the correct approach when the goal is predictive
performance at scale on tasks where structure is implicit and data is
abundant. Field-theoretic dialectical frameworks may be the correct
approach when the goal is \textbf{structural understanding of attractor
dynamics in low-data or interpretability-critical settings} (scientific
hypothesis generation, safety-critical decision systems, educational /
explanatory applications). Neither paradigm subsumes the other.

\textbf{The long-term question} is whether the two paradigms can be
composed --- statistical deep learning providing the upstream embedding
(as in our \texttt{F\_dense} construction with BGE-M3) and the
field-theoretic framework providing the downstream structural analysis.
Our empirical results suggest this composition is \textbf{non-trivial}:
BGE's statistical structure violates Laplace's mean-zero source
assumption by 17σ (§4.9), a quantitative obstruction to naïve
composition. Addressing this obstruction is OP2 sub-problem (a) (§5.2.2)
and Block V Phase A-1 (§5.3). The outcome will substantively inform how
the two paradigms can interoperate.

\hypertarget{implications-for-the-dialectical-materialist-tradition}{%
\subsection{6.5 Implications for the Dialectical-Materialist
Tradition}\label{implications-for-the-dialectical-materialist-tradition}}

This paper is an explicit attempt to \textbf{re-activate dialectical
materialism as a working method} for 21st-century problems in artificial
intelligence, beyond its 20th-century applications to political economy
and social theory. We take seriously three commitments from that
tradition:

\begin{enumerate}
\def\labelenumi{\arabic{enumi}.}
\tightlist
\item
  That \textbf{practice is the criterion of truth} (On Practice, Mao
  1937). Our spawn-agent review pipeline, our mode-tagging discipline,
  and our explicit OP-signposting are all practical instruments, not
  rhetorical positions; they are the framework's working-out of the
  practice criterion at the level of paper production
\item
  That \textbf{contradictions internal to a system drive its motion} (On
  Contradiction, Mao 1937). We have been explicit about the framework's
  internal contradictions (Axiom 4 vs.~BGE, Axiom 3 vs.~Langevin, Axiom
  7 vs.~fixed V) rather than resolving them by fiat; their persistence
  is what the framework needs to continue developing
\item
  That \textbf{quantitative changes can produce qualitative transitions}
  (Engels, posthumously 1925). Stage A1.4 (§4.9) is the framework's own
  direct experimental instance of this principle; the ``failure'' of
  A1.4 --- its 52.5\% clamp residency --- is not an engineering setback
  but a quality-transition observation that refines OP2 into its
  three-fold form
\end{enumerate}

We do not claim to speak for the tradition beyond our specific use of
it. Other dialectical-materialist thinkers have developed different
methodological emphases (structuralism, historical sociology, critical
theory); our use is technical and instrument-focused, not programmatic.

\textbf{On the broader project}. We hold the position, defensible
against critical examination, that dialectical materialism is a
methodological tradition with substantial unresolved potential for
application to 21st-century problems --- specifically, that the
tradition's analytical tools for understanding how systems evolve
through internal contradiction may be more applicable to AI and complex
systems than to the primarily social-scientific applications that
dominated 20th-century practice. This paper is an experiment in testing
that position; the results are, as the paper shows, partial but
substantive.

We do not expect this framing to be universally welcomed. We include it
to be honest about the paper's motivation --- and because the reader who
knows the motivation can more accurately calibrate what to believe and
what to question.

\hypertarget{concluding-statement}{%
\subsection{6.6 Concluding Statement}\label{concluding-statement}}

A framework whose research program is entirely in its open problems is
not an evasion of responsibility. It is the honest representation of a
research stage in which the most important claims are \textbf{what we
have not proven} rather than \textbf{what we have proven}. We have:

\begin{itemize}
\tightlist
\item
  Proposed a categorical structure with explicit conjectural limits
\item
  Reported empirical findings with explicit finite-lattice caveats
\item
  Identified two open problems with concrete experimental programs
\item
  Established a three-stage publication plan that separates mechanism
  from integration
\end{itemize}

The next paper in the series (the IPM mechanism paper, Stage 2, target
2026-07--08 month) will report Block V Phase A outcomes. The capstone
paper (Stage 3, target 2027 Q1--Q2) will integrate across Phases A, B,
and C with whatever state the research has reached. Neither is promised.

The framework's ultimate target is \textbf{not a benchmark victory}. The
target is \textbf{the re-activation of dialectical materialism as a
working methodology for 21st-century productive forces} --- a target
measured in years to decades, not quarters, and one whose success
criterion is the emergence of a community that finds the framework
useful, not any single empirical result. We believe the work reported
here is a coherent first anchor toward that target; we do not believe it
is the target. The distinction matters.

\begin{center}\rule{0.5\linewidth}{0.5pt}\end{center}

\hypertarget{references}{%
\section{References}\label{references}}

\textbf{Primary references cited in text} (full bibliographic details in
\texttt{references.bib} in the companion repository; reproducibility
scripts in \texttt{paper/reproduce/}):

\begin{itemize}
\tightlist
\item
  Lawvere, F. W. (1969). \emph{Adjointness in foundations}. Dialectica,
  23:281--296.
\item
  Mao Zedong (1937). \emph{On Contradiction} (矛盾论).
\item
  Mao Zedong (1937). \emph{On Practice} (实践论).
\item
  Lenin, V. I. (1909). \emph{Materialism and Empirio-Criticism}.
\item
  Engels, F. (posthumously 1925). \emph{Dialectics of Nature}.
\item
  Giry, M. (1982). \emph{A categorical approach to probability theory}.
  In Categorical Aspects of Topology and Analysis, LNM 915, Springer.
\item
  Jacobs, B. (2010). \emph{Convexity, duality and effects}. In
  Theoretical Computer Science, Springer.
\item
  Adams, H., et al.~(2021). \emph{giotto-tda: A Topological Data
  Analysis Toolkit for Machine Learning and Data Exploration}. JMLR.
\item
  Bossy, M., \& Talay, D. (2005). \emph{Reflected Langevin dynamics and
  related schemes}. Stochastic Analysis and Applications.
\item
  Anderson, P. W. (1972). \emph{More Is Different}. Science
  177:393--396.
\item
  Kramers, H. A. (1940). \emph{Brownian motion in a field of force and
  the diffusion model of chemical reactions}. Physica 7:284--304.
\end{itemize}

\textbf{Supporting methodology documents} (in this repository):

\begin{itemize}
\tightlist
\item
  \texttt{FINAL\_REPORT.md} --- Full exp017 experimental archive (Block
  I--IV + Stage C/A1)
\item
  \texttt{block4\_5/a1\_4/VERDICT.md} --- A1.4 shifted-potential verdict
  (Phase 2/2.5)
\item
  \texttt{BLOCK\_V\_DESIGN.md} --- Phase A/B/C experimental roadmap
\item
  \texttt{REVIEW\_*.md} files --- Spawn-agent review records (P1--P5 +
  Phase 2.5 + §4 + verdicts)
\end{itemize}

\begin{center}\rule{0.5\linewidth}{0.5pt}\end{center}

\end{document}
